% =====================================================================
%  EXTRACTO PARA REVISIÓN — PARTE I (TRONCO COMÚN)
%
%  Documento aparte para mandar a revisar solo lo que está escrito.
%  NO duplica contenido: usa los mismos ficheros de capítulo que main.tex.
%
%  Compilar con:  exportar_revision.bat
%
%  CÓMO FUNCIONA:
%  Los capítulos de la parte I se incluyen normalmente. De los demás solo se
%  leen las etiquetas de su .aux con \includeref, para que las referencias
%  cruzadas a la parte no escrita (cap. 53, 26, 13...) sigan saliendo bien en
%  vez de «??».  \includeonly no sirve aquí: arrastra al índice las entradas
%  de los capítulos excluidos y el índice acabaría listando los 60.
%
%  REQUIERE haber compilado main.tex antes: de ahí salen esos .aux.
%  El .bat se encarga del orden.
% =====================================================================
\documentclass[11pt,a4paper,oneside,openany]{book}

% =====================================================================
%  PREÁMBULO — Curso de Ingeniería FUSAL
%  Todo lo que configura el documento vive aquí. main.tex se queda limpio.
% =====================================================================

% ---- Idioma y codificación -------------------------------------------
\usepackage[utf8]{inputenc}
\usepackage[T1]{fontenc}
\usepackage[spanish,es-noshorthands]{babel}   % es-noshorthands evita choques con listings/tikz/siunitx
\usepackage{csquotes}

% ---- Tipografía ------------------------------------------------------
\usepackage{amsmath,amssymb,mathtools}
\usepackage{newtxtext}          % Texto tipo Times
\usepackage{newtxmath}          % Matemáticas a juego
\usepackage{microtype}
% Para cambiar de fuente: sustituye las dos líneas newtx por
%   \usepackage{XCharter}\usepackage[charter]{mathdesign}   (Charter, más legible en pantalla)
%   \usepackage{palatino}                                    (Palatino)

% ---- Geometría de página --------------------------------------------
\usepackage[a4paper,left=3cm,right=2.5cm,top=3cm,bottom=3cm,headheight=15pt]{geometry}
\usepackage{setspace}
\onehalfspacing
\setlength{\parskip}{0.4em}

% ---- Colores del equipo ---------------------------------------------
\usepackage{xcolor}
\definecolor{fusalPrim}{HTML}{1F3864}   % azul oscuro  <- cambia aquí el color corporativo
\definecolor{fusalSec}{HTML}{C00000}    % rojo acento
\definecolor{fusalGris}{HTML}{6E6E6E}
\definecolor{fusalClaro}{HTML}{EEF2F8}

% ---- Figuras y tablas ------------------------------------------------
\usepackage{graphicx}
\graphicspath{{Imagenes/}{Imagenes/dinamica/}{Imagenes/suspension/}%
              {Imagenes/chasis/}{Imagenes/aero/}{Imagenes/refrigeracion/}%
              {Imagenes/transmision/}{Imagenes/transversal/}}
\usepackage{float}
\usepackage{subcaption}
\usepackage[font=small,labelfont=bf,labelsep=period]{caption}
\usepackage{booktabs,tabularx,longtable,multirow,array}
\usepackage{rotating}
\usepackage{pdfpages}
\usepackage{enumitem}
\setlist{itemsep=2pt,topsep=4pt}

% ---- Unidades --------------------------------------------------------
\usepackage{siunitx}
\sisetup{
  output-decimal-marker = {,},   % coma decimal (convenio español)
  per-mode              = symbol,
  separate-uncertainty  = true,
  range-phrase          = {--},
  range-units           = single,
  detect-weight         = true
}
\DeclareSIUnit{\rpm}{rpm}
\DeclareSIUnit{\bar}{bar}
\DeclareSIUnit{\gacc}{g}

% ---- Encabezados y pies ---------------------------------------------
\usepackage{fancyhdr}
\pagestyle{fancy}
\fancyhf{}
\fancyhead[L]{\small\textcolor{fusalGris}{\nouppercase{\leftmark}}}
\fancyhead[R]{\small\textcolor{fusalGris}{\thepage}}
\renewcommand{\headrulewidth}{0.4pt}
\fancypagestyle{plain}{\fancyhf{}\fancyfoot[C]{\small\thepage}\renewcommand{\headrulewidth}{0pt}}

% ---- Estilo de títulos ----------------------------------------------
\usepackage{titlesec}
\titleformat{\chapter}[hang]
  {\normalfont\sffamily\huge\bfseries\color{fusalPrim}}
  {\thechapter\hspace{10pt}\textcolor{fusalGris}{\textbar}\hspace{10pt}}{0pt}{\huge}
\titlespacing*{\chapter}{0pt}{-15pt}{28pt}
\titleformat{\section}
  {\normalfont\sffamily\Large\bfseries\color{fusalPrim}}{\thesection}{1em}{}
\titleformat{\subsection}
  {\normalfont\sffamily\large\bfseries\color{fusalGris}}{\thesubsection}{1em}{}

% ---- Índices ---------------------------------------------------------
\usepackage{makeidx}
\makeindex
\setcounter{tocdepth}{1}      % en el índice general, solo hasta \section
\setcounter{secnumdepth}{3}

% ---- Código fuente (carpeta Codigo/) ---------------------------------
\usepackage{listings}
\definecolor{codekw}{HTML}{0033B3}
\definecolor{codecmt}{HTML}{5C8A4A}
\definecolor{codestr}{HTML}{A31515}
\definecolor{codebg}{HTML}{F7F7F7}
\definecolor{codenum}{HTML}{999999}
\lstdefinestyle{cursocode}{
    basicstyle=\ttfamily\scriptsize,
    keywordstyle=\color{codekw}\bfseries,
    commentstyle=\color{codecmt}\itshape,
    stringstyle=\color{codestr},
    numbers=left, numberstyle=\tiny\color{codenum}, numbersep=7pt,
    backgroundcolor=\color{codebg},
    frame=single, rulecolor=\color{codenum!60},
    breaklines=true, breakatwhitespace=false, showstringspaces=false,
    tabsize=4, columns=fullflexible, keepspaces=true,
    extendedchars=true, inputencoding=utf8, literate=%
        {á}{{\'a}}1 {é}{{\'e}}1 {í}{{\'i}}1 {ó}{{\'o}}1 {ú}{{\'u}}1
        {ñ}{{\~n}}1 {Á}{{\'A}}1 {É}{{\'E}}1 {Í}{{\'I}}1 {Ó}{{\'O}}1
        {Ú}{{\'U}}1 {Ñ}{{\~N}}1 {¿}{{?`}}1 {¡}{{!`}}1 {°}{{\textdegree}}1
        {²}{{\textsuperscript{2}}}1 {³}{{\textsuperscript{3}}}1
        {µ}{{\textmu}}1 {·}{{\textperiodcentered}}1 {×}{{\texttimes}}1
        {–}{{--}}1 {—}{{---}}1 {•}{{\textbullet}}1
}
\lstset{style=cursocode,inputpath=Codigo/}
\renewcommand{\lstlistingname}{Código}
\renewcommand{\lstlistlistingname}{Índice de códigos}

% ---- Dibujo ----------------------------------------------------------
\usepackage{tikz}
\usetikzlibrary{arrows.meta,calc,positioning,patterns,decorations.pathmorphing,decorations.pathreplacing}
\usepackage{pgfplots}
\usepgfplotslibrary{groupplots,fillbetween}
\pgfplotsset{compat=1.18}

% Estilo común de todas las gráficas del libro. Úsalo siempre:
%   \begin{axis}[cursoplot, xlabel=..., ylabel=...]   una gráfica a ancho de página
%   \begin{axis}[cursopanel, ...]                     media página (dos en paralelo)
\pgfplotsset{/pgf/number format/.cd, use comma, 1000 sep={}}
\pgfplotsset{
  cursoplot/.style={
    width=0.92\linewidth, height=6.2cm,
    scale only axis, axis lines=left,
    label style={font=\sffamily\small},
    tick label style={font=\sffamily\scriptsize},
    title style={font=\sffamily\small\bfseries,color=fusalPrim},
    legend style={font=\sffamily\scriptsize,draw=fusalGris!50,fill=white,
                  fill opacity=0.88,text opacity=1,rounded corners=1pt},
    legend cell align=left,
    grid=both, major grid style={draw=fusalGris!25},
    minor grid style={draw=fusalGris!12},
    scaled ticks=false, tick align=outside,
    every axis plot/.append style={line width=0.9pt},
    clip mode=individual},
  % Panel para dos gráficas en paralelo. Aquí «width» es el ancho TOTAL del
  % dibujo (rótulos incluidos), por eso lleva scale only axis=false.
  cursopanel/.style={cursoplot, scale only axis=false,
    width=0.47\linewidth, height=6.0cm},
  cursotercio/.style={cursoplot, scale only axis=false,
    width=0.32\linewidth, height=5.4cm,
    label style={font=\sffamily\scriptsize},
    tick label style={font=\sffamily\tiny},
    title style={font=\sffamily\scriptsize\bfseries,color=fusalPrim},
    legend style={font=\sffamily\tiny,draw=fusalGris!50,fill=white,
                  fill opacity=0.88,text opacity=1,rounded corners=1pt}},
}
% Magic Formula de Pacejka, para dibujar curvas de neumático (capítulos 8 y 9):
%   mf(x,B,C,D,E) = D sin( C arctan( (1-E) Bx + E arctan(Bx) ) )
% OJO con los grados: en pgfmath atan() devuelve grados y sin() los espera, por
% eso la arcotangente interior lleva rad(). Va aquí y no en el capítulo porque
% declare function es global y no se puede declarar dos veces.
\tikzset{
  declare function={
    mf(\x,\B,\C,\D,\E) = \D*sin(\C*atan((1-\E)*\B*\x + \E*rad(atan(\B*\x))));
  }
}

% Colores de curva. El «cycle list» automático de pgfplots no es fiable en todas
% las instalaciones, así que cada \addplot lleva su estilo explícito:
%   \addplot[curvaA] ... \addplot[curvaB] ...
\tikzset{
  curvaA/.style={fusalPrim,thick},
  curvaB/.style={fusalSec,thick},
  curvaC/.style={teal!65!black,thick},
  curvaD/.style={orange!85!black,thick},
  curvaE/.style={violet!65!black,thick},
  curvaAux/.style={fusalGris,thick,densely dashed},
  cursonota/.style={font=\sffamily\scriptsize,align=center,text=fusalGris},
  cursomarca/.style={fusalSec,densely dashed,line width=0.6pt},
}

% ---- Bibliografía ----------------------------------------------------
\usepackage[style=numeric-comp,backend=biber,sorting=none,maxbibnames=99]{biblatex}
\addbibresource{bibliografia.bib}

% ---- Notas de trabajo (modo borrador) --------------------------------
% Con \borradortrue (en main.tex) se muestran los \todo y los avisos.
\newif\ifborrador
\usepackage[colorinlistoftodos,prependcaption,textsize=footnotesize,spanish]{todonotes}
\setlength{\marginparwidth}{2cm}   % silencia el aviso de todonotes (solo usamos todos en línea)
\newcommand{\pendiente}[1]{\ifborrador\todo[inline,color=yellow!40]{#1}\fi}
\newcommand{\falta}[1]{\ifborrador\todo[inline,color=red!25]{FALTA: #1}\fi}
\newcommand{\figpendiente}[1]{\ifborrador\missingfigure{#1}\fi}

% ---- Hiperenlaces (siempre lo último salvo cleveref) -----------------
\usepackage[hidelinks,bookmarksnumbered]{hyperref}
\hypersetup{
  colorlinks = true,
  linkcolor  = fusalPrim,
  citecolor  = fusalSec,
  urlcolor   = fusalPrim,
  pdftitle   = {Curso de Ingeniería FUSAL — Mechanical},
  pdfauthor  = {Formula Student USAL}
}
\usepackage[nameinlink]{cleveref}
\crefname{figure}{figura}{figuras}
\crefname{table}{tabla}{tablas}
\crefname{equation}{ecuación}{ecuaciones}
\crefname{chapter}{capítulo}{capítulos}
\crefname{section}{apartado}{apartados}

% ---- Entornos propios y macros --------------------------------------
% =====================================================================
%  ENTORNOS PROPIOS
%  Los seis cuadros que estructuran cada capítulo. Úsalos siempre igual:
%  el lector debe poder hojear el libro buscando solo los recuadros.
% =====================================================================
\usepackage[most]{tcolorbox}
\tcbuselibrary{breakable,skins}

% --- Plantilla base ---------------------------------------------------
\tcbset{
  cursobase/.style={
    breakable, enhanced, sharp corners=downhill,
    boxrule=0pt, leftrule=3pt, toprule=0pt, bottomrule=0pt, rightrule=0pt,
    left=8pt, right=8pt, top=6pt, bottom=6pt,
    fonttitle=\sffamily\bfseries\small,
    coltitle=black, attach title to upper=\par\medskip
  }
}

% --- 1. Idea clave ----------------------------------------------------
\newtcolorbox{clave}[1][]{cursobase,
  colback=fusalClaro, colframe=fusalPrim,
  title={IDEA CLAVE}, #1}

% --- 2. Caso real del equipo -----------------------------------------
\newtcolorbox{caso}[1][]{cursobase,
  colback=fusalPrim!5, colframe=fusalPrim!70,
  title={CASO FUSAL\ifx\relax#1\relax\else\ --- #1\fi}}

% --- 3. Referencia al reglamento -------------------------------------
\newtcolorbox{regla}[1][]{cursobase,
  colback=gray!8, colframe=fusalGris,
  title={DEL REGLAMENTO\ifx\relax#1\relax\else\ --- #1\fi}}

% --- 4. Errores frecuentes -------------------------------------------
\newtcolorbox{ojo}[1][]{cursobase,
  colback=fusalSec!5, colframe=fusalSec,
  title={OJO --- ERRORES FRECUENTES}, #1}

% --- 5. Entregable de la sesión --------------------------------------
\newtcolorbox{entregable}[1][]{cursobase,
  colback=green!5, colframe=green!45!black,
  title={ENTREGABLE}, #1}

% --- 6. Preguntas de los jueces --------------------------------------
\newtcolorbox{juez}[1][]{cursobase,
  colback=orange!6, colframe=orange!70!black,
  title={LO QUE TE VA A PREGUNTAR UN JUEZ}, #1}

% --- Objetivos de aprendizaje (cabecera de cada capítulo) -------------
\newenvironment{objetivos}
  {\begin{tcolorbox}[cursobase, colback=white, colframe=fusalPrim,
      title={AL TERMINAR ESTE CAPÍTULO DEBES SABER}]
   \begin{itemize}[leftmargin=*,label=\textcolor{fusalPrim}{\textbullet}]}
  {\end{itemize}\end{tcolorbox}}

% --- Ficha de sesión (para impartir el capítulo como clase) -----------
% Uso: \fichasesion{Duración}{Requisitos previos}{Material necesario}
\newcommand{\fichasesion}[3]{%
  \begin{tcolorbox}[cursobase, colback=gray!5, colframe=fusalGris,
      title={FICHA DE SESIÓN}]
  \begin{description}[leftmargin=!,labelwidth=2.6cm,font=\sffamily\bfseries\small]
    \item[Duración] #1
    \item[Requisitos] #2
    \item[Material] #3
  \end{description}
  \end{tcolorbox}}

% =====================================================================
%  MACROS
%  Símbolos que se repiten en todo el libro. Escribe SIEMPRE con macro,
%  nunca a mano: si un día cambias el convenio, cambias solo esta hoja.
% =====================================================================

% --- Generales --------------------------------------------------------
\usepackage{xspace}
\newcommand{\FUSAL}{\textsc{FUSAL}\xspace}
\newcommand{\etal}{\textit{et al.}\xspace}
\newcommand{\dif}{\mathrm{d}}
\newcommand{\deriv}[2]{\frac{\dif #1}{\dif #2}}
\newcommand{\pderiv}[2]{\frac{\partial #1}{\partial #2}}
\newcommand{\vect}[1]{\mathbf{#1}}
\newcommand{\prom}[1]{\overline{#1}}
% Texto guía dentro de los recuadros: qué hay que escribir ahí. Bórralo al redactar.
\newcommand{\redactar}[1]{\textcolor{fusalGris}{\itshape #1}}

% --- Dinámica vehicular ----------------------------------------------
\newcommand{\Fz}{F_z}                       % carga vertical
\newcommand{\Fy}{F_y}                       % fuerza lateral
\newcommand{\Fx}{F_x}                       % fuerza longitudinal
\newcommand{\Mz}{M_z}                       % par autoalineante
\newcommand{\slipa}{\alpha}                 % ángulo de deriva
\newcommand{\slipr}{\kappa}                 % deslizamiento longitudinal
\newcommand{\camber}{\gamma}                % caída
\newcommand{\muy}{\mu_y}
\newcommand{\mux}{\mu_x}
\newcommand{\ay}{a_y}
\newcommand{\ax}{a_x}
\newcommand{\CdG}{\mathrm{CdG}}             % centro de gravedad
\newcommand{\hcg}{h_{\CdG}}
\newcommand{\LLTD}{\mathrm{LLTD}}
\newcommand{\Kus}{K_{us}}                   % gradiente de subviraje
\newcommand{\batalla}{L}
\newcommand{\via}{t}
\newcommand{\Kroll}{K_{\phi}}               % rigidez a balanceo

% --- Neumático y su modelado (capítulos 8 y 9) ------------------------
\newcommand{\Calpha}{C_{\slipa}}            % rigidez de deriva
\newcommand{\Ckappa}{C_{\slipr}}            % rigidez de deslizamiento
\newcommand{\tneu}{t_p}                     % avance neumático
\newcommand{\pneu}{p_i}                     % presión de inflado
\newcommand{\Tneu}{T_n}                     % temperatura del neumático
\newcommand{\Fzo}{F_{z0}}                   % carga vertical de referencia
\newcommand{\dfz}{\mathrm{d}f_z}            % variación normalizada de carga
\newcommand{\lamu}{\lambda_{\mu}}           % factor de escala de agarre
% Coeficientes de la Magic Formula. OJO: \Cp ya es el coeficiente de
% presión aerodinámico, por eso estos llevan prefijo MF.
\newcommand{\MFB}{B}                        % factor de rigidez
\newcommand{\MFC}{C}                        % factor de forma
\newcommand{\MFD}{D}                        % valor de pico
\newcommand{\MFE}{E}                        % factor de curvatura
\newcommand{\MFSh}{S_h}                     % desplazamiento horizontal
\newcommand{\MFSv}{S_v}                     % desplazamiento vertical

% --- Suspensión -------------------------------------------------------
\newcommand{\MR}{\mathrm{MR}}               % relación de instalación
\newcommand{\Kw}{K_w}                       % rigidez en rueda
\newcommand{\Ks}{K_s}                       % rigidez del muelle
\newcommand{\fr}{f_r}                       % frecuencia de suspensión

% --- Aerodinámica -----------------------------------------------------
\newcommand{\Cl}{C_L}
\newcommand{\Cd}{C_D}
\newcommand{\Cp}{C_p}
\newcommand{\ClA}{C_L A}
\newcommand{\CdA}{C_D A}
\newcommand{\qdin}{q_{\infty}}              % presión dinámica
\newcommand{\Vinf}{V_{\infty}}
\newcommand{\Rey}{\mathit{Re}}
\newcommand{\CoP}{\mathrm{CoP}}             % centro de presiones

% --- Transferencia de calor ------------------------------------------
\newcommand{\UA}{UA}
\newcommand{\NTU}{\mathit{NTU}}
\newcommand{\efec}{\varepsilon}             % eficacia del intercambiador
\newcommand{\Nus}{\mathit{Nu}}
\renewcommand{\Pr}{\mathit{Pr}}   % redefine el \Pr de LaTeX (probabilidad), aquí es Prandtl
\newcommand{\Stan}{\mathit{St}}
\newcommand{\jcol}{j}                       % factor de Colburn
\newcommand{\ffan}{f}                       % factor de fricción
\newcommand{\mdot}{\dot{m}}
\newcommand{\Qdot}{\dot{Q}}
\newcommand{\Dh}{D_h}                       % diámetro hidráulico
\newcommand{\LMTD}{\Delta T_{\mathrm{ml}}}

% --- Transmisión ------------------------------------------------------
\newcommand{\itot}{i_{\mathrm{tot}}}        % relación total
\newcommand{\Tmot}{T_{\mathrm{mot}}}
\newcommand{\rdin}{r_{\mathrm{din}}}        % radio dinámico

% =====================================================================
%  ESTILOS DE DIAGRAMA DE FLUJO
%  Usados en los organigramas del libro. Convenio:
%    proceso    rectángulo      una acción
%    decision   rombo           una pregunta con salidas Sí/No
%    terminal   elipse          principio o final del flujo
%    nota       recuadro gris   detalle de lo que incluye un paso
%    conector   triángulo       enlace a un subproceso desarrollado aparte
%    propuesto  trazo discontinuo rojo: añadido, todavía no acordado
% =====================================================================
\usetikzlibrary{shapes.geometric}

\tikzset{
  proceso/.style={
    draw=fusalPrim!75, fill=fusalPrim!6, rounded corners=2pt,
    align=center, inner sep=4pt, minimum height=8mm, text width=26mm,
    font=\sffamily\scriptsize},
  decision/.style={
    draw=fusalPrim!75, fill=white, diamond, aspect=1.7,
    align=center, inner sep=0pt, text width=19mm,
    font=\sffamily\scriptsize},
  terminal/.style={
    draw=fusalPrim, fill=fusalPrim!15, ellipse,
    align=center, inner sep=3pt, text width=28mm,
    font=\sffamily\scriptsize\bfseries},
  nota/.style={
    draw=fusalGris, fill=gray!7, align=left, inner sep=4pt,
    text width=32mm, font=\sffamily\tiny},
  conector/.style={
    draw=fusalSec, fill=fusalSec!12, regular polygon,
    regular polygon sides=3, inner sep=0pt, minimum size=7mm,
    font=\sffamily\tiny\bfseries, text=fusalSec},
  propuesto/.style={
    draw=fusalSec, fill=fusalSec!5, dashed, dash pattern=on 2pt off 1.5pt},
  flujo/.style={-{Latex[length=2mm,width=1.6mm]}, draw=fusalPrim!85, semithick},
  flujoprop/.style={-{Latex[length=2mm,width=1.6mm]}, draw=fusalSec,
    dashed, dash pattern=on 2pt off 1.5pt, semithick},
  aviso/.style={-{Latex[length=1.6mm,width=1.3mm]}, draw=fusalGris},
  etiq/.style={font=\sffamily\tiny, inner sep=1.5pt, fill=white}
}



% ---- Datos del documento --------------------------------------------
\newcommand{\temporada}{2025--2026}
\newcommand{\version}{0.1 --- extracto para revisión}

% Modo borrador ACTIVADO a propósito: los revisores deben ver qué falta.
\borradortrue

% ---------------------------------------------------------------------
%  \includeref{fichero}
%  Lee las etiquetas del .aux de un capítulo que NO se imprime, para que
%  \ref{ch:evento_costes} siga dando 53 sin que el capítulo 53 aparezca en
%  este extracto.
%
%  Tiene que llamarse EN EL PREÁMBULO, por dos motivos:
%    - dentro del documento, «@» ya no es letra y el .aux no se puede leer;
%    - en el preámbulo los índices aún no están abiertos, así que las
%      entradas de índice de esos capítulos se descartan solas.
% ---------------------------------------------------------------------
\makeatletter
\newcommand{\includeref}[1]{\begingroup\makeatletter\@input{#1.aux}\endgroup}
\makeatother

% ---------------------------------------------------------------------
%  Partes II a VIII y anexos: NO se imprimen. Solo se leen sus etiquetas.
%  Al añadir capítulos nuevos al libro, añádelos también a esta lista.
% ---------------------------------------------------------------------
\includeref{Capitulos/02_Dinamica/08_neumatico}
\includeref{Capitulos/02_Dinamica/09_modelado_neumatico}
\includeref{Capitulos/02_Dinamica/10_regimen_estacionario}
\includeref{Capitulos/02_Dinamica/11_diagramas_ymd}
\includeref{Capitulos/02_Dinamica/12_regimen_transitorio}
\includeref{Capitulos/02_Dinamica/13_simulacion_vuelta}
\includeref{Capitulos/02_Dinamica/14_puesta_a_punto}
\includeref{Capitulos/02_Dinamica/15_piloto_ensayos_pista}

\includeref{Capitulos/03_Suspension/16_cinematica_suspension}
\includeref{Capitulos/03_Suspension/17_puntos_duros}
\includeref{Capitulos/03_Suspension/18_muelles_amortiguadores}
\includeref{Capitulos/03_Suspension/19_casos_carga}
\includeref{Capitulos/03_Suspension/20_componentes_suspension}
\includeref{Capitulos/03_Suspension/21_direccion}
\includeref{Capitulos/03_Suspension/22_frenos}
\includeref{Capitulos/03_Suspension/23_montaje_alineacion}

\includeref{Capitulos/04_Chasis/24_chasis_reglamento}
\includeref{Capitulos/04_Chasis/25_tubular_vs_monocasco}
\includeref{Capitulos/04_Chasis/26_rigidez_torsional}
\includeref{Capitulos/04_Chasis/27_caminos_carga_nodos}
\includeref{Capitulos/04_Chasis/28_materiales_compuestos}
\includeref{Capitulos/04_Chasis/29_fabricacion_tubular}
\includeref{Capitulos/04_Chasis/30_estructuras_impacto}
\includeref{Capitulos/04_Chasis/31_ergonomia_packaging}

\includeref{Capitulos/05_Aerodinamica/32_fundamentos_aerodinamica}
\includeref{Capitulos/05_Aerodinamica/33_merece_la_pena_aero}
\includeref{Capitulos/05_Aerodinamica/34_elementos_aerodinamicos}
\includeref{Capitulos/05_Aerodinamica/35_cfd_montaje_caso}
\includeref{Capitulos/05_Aerodinamica/36_cfd_validacion}
\includeref{Capitulos/05_Aerodinamica/37_ensayos_pista_aero}
\includeref{Capitulos/05_Aerodinamica/38_fabricacion_aero}
\includeref{Capitulos/05_Aerodinamica/39_interaccion_aero_suspension}

\includeref{Capitulos/06_Refrigeracion/40_balance_termico}
\includeref{Capitulos/06_Refrigeracion/41_fundamentos_intercambiadores}
\includeref{Capitulos/06_Refrigeracion/42_superficies_compactas}
\includeref{Capitulos/06_Refrigeracion/43_lado_aire}
\includeref{Capitulos/06_Refrigeracion/44_lado_liquido}
\includeref{Capitulos/06_Refrigeracion/45_validacion_termica}

\includeref{Capitulos/07_Transmision/46_arquitectura_transmision}
\includeref{Capitulos/07_Transmision/47_relacion_transmision}
\includeref{Capitulos/07_Transmision/48_elementos_transmision}
\includeref{Capitulos/07_Transmision/49_diferencial}
\includeref{Capitulos/07_Transmision/50_palieres_juntas}
\includeref{Capitulos/07_Transmision/51_soportes_seguridad}
\includeref{Capitulos/07_Transmision/52_eficiencia_energetica}

\includeref{Capitulos/08_Transversal/53_evento_costes}
\includeref{Capitulos/08_Transversal/54_coste_real_compras}
\includeref{Capitulos/08_Transversal/55_fabricacion_i}
\includeref{Capitulos/08_Transversal/56_fabricacion_ii}
\includeref{Capitulos/08_Transversal/57_dfm_dfa}
\includeref{Capitulos/08_Transversal/58_planificacion_fiabilidad}
\includeref{Capitulos/08_Transversal/59_design_event}
\includeref{Capitulos/08_Transversal/60_documentacion_traspaso}

\includeref{Capitulos/99_Anexos/A_plantillas}
\includeref{Capitulos/99_Anexos/B_checklists}
\includeref{Capitulos/99_Anexos/C_planos}
\includeref{Capitulos/99_Anexos/D_codigo}

\begin{document}

% =====================================================================
\frontmatter
% =====================================================================
% =====================================================================
%  PORTADA
%  Sustituye el rectángulo por el logo del equipo cuando lo tengas:
%  \includegraphics[width=6cm]{Imagenes/logo_fusal.pdf}
% =====================================================================
\begin{titlepage}
\centering
\thispagestyle{empty}

\vspace*{1.5cm}

{\sffamily\Large\color{fusalGris} FORMULA STUDENT · UNIVERSIDAD DE SALAMANCA\par}

\vspace{1.2cm}

% --- Sustituir por el logo ---
\begin{tikzpicture}
  \draw[fusalGris!40,line width=0.8pt] (0,0) rectangle (6,2.2);
  \node[fusalGris] at (3,1.1) {\sffamily\small LOGO FUSAL};
\end{tikzpicture}

\vspace{1.6cm}

{\sffamily\bfseries\Huge\color{fusalPrim} Curso de Ingeniería \FUSAL\par}

\vspace{0.6cm}

{\sffamily\LARGE\color{fusalGris} Departamento de Mecánica\par}

\vspace{0.4cm}
{\color{fusalSec}\rule{9cm}{2pt}}
\vspace{0.8cm}

{\large\itshape Chasis · Aerodinámica · Refrigeración · Dinámica vehicular\\
Suspensión · Transmisión · Costes y fabricación\par}

\vfill

{\sffamily\large Manual de formación y transferencia de conocimiento\par}
\vspace{0.3cm}
{\sffamily\normalsize\color{fusalGris} Temporada \temporada\ \textbar\ Versión \version\par}

\vspace{1cm}
{\small\color{fusalGris}\today\par}

\end{titlepage}

% --- Página de créditos ---------------------------------------------
\thispagestyle{empty}
\null\vfill
\noindent
\begin{minipage}{\textwidth}
\small
\textbf{Curso de Ingeniería \FUSAL{} --- Departamento de Mecánica}\\[0.4em]
Documento interno del equipo Formula Student de la Universidad de Salamanca.\\[0.8em]

\textbf{Cómo citar este documento:}\\
Equipo \FUSAL. \textit{Curso de Ingeniería \FUSAL: Mechanical}. Versión \version, temporada \temporada.\\[0.8em]

\textbf{Autores y revisores por capítulo:} véase la tabla de responsables en el
apartado \textit{Cómo usar este libro}.\\[0.8em]

\textbf{Aviso.} Los valores numéricos, reglas y procedimientos recogidos aquí
corresponden a la temporada indicada. \emph{El reglamento cambia todos los años}:
antes de aplicar cualquier requisito normativo de este libro, contrástalo con la
versión vigente del reglamento.\\[0.8em]

\textbf{Erratas y mejoras:} anótalas en el registro del equipo. Este libro se
actualiza cada temporada; si detectas algo mal y no lo escribes, se pierde.
\end{minipage}
\vspace{2cm}
\clearpage


% !TeX root = ../../CursoFUSAL_Bloque0_revision.tex
% =====================================================================
%  Nota de revisión
%  SOLO aparece en el extracto que se manda a revisar (no en main.tex)
% =====================================================================
\chapter*{Cómo revisar este documento}
\addcontentsline{toc}{chapter}{Cómo revisar este documento}
\markboth{Cómo revisar este documento}{Cómo revisar este documento}

\section*{Qué es esto}

Este documento es un \textbf{extracto} del \textit{Curso de Ingeniería \FUSAL}:
contiene únicamente la parte I, el tronco común, que es la que se lee entera todo
el equipo antes de tocar el CAD. El libro completo tiene ocho partes y sesenta
capítulos; el resto todavía está sin redactar.

La numeración de capítulos y las referencias cruzadas son las del libro completo.
Por eso verás llamadas a capítulos que aquí no aparecen (por ejemplo, «capítulo
\ref{ch:evento_costes}»): son correctas, pero remiten a partes todavía no escritas.

\section*{Estado de cada capítulo}

\begin{center}
\begin{tabularx}{\textwidth}{@{}clX@{}}
\toprule
\textbf{Cap.} & \textbf{Título} & \textbf{Estado} \\
\midrule
\ref{ch:formula_student} & Formula Student: el juego y sus reglas & Redactado \\
\ref{ch:coche_como_sistema} & El coche como sistema & Redactado, breve a propósito \\
\ref{ch:proceso_diseno} & Proceso de diseño & \textbf{Solo los diagramas.} El texto está pendiente \\
\ref{ch:herramientas_cad_datos} & Herramientas: CAD, datos y versiones & Redactado \\
\ref{ch:mecanica_aplicada} & Fundamentos de mecánica aplicada & Redactado \\
\ref{ch:elementos_finitos} & Análisis por elementos finitos & Redactado \\
\ref{ch:ensayo_adquisicion_datos} & Ensayo y adquisición de datos & \textbf{Vacío a propósito}, se explica por qué \\
\bottomrule
\end{tabularx}
\end{center}

\section*{Sobre qué queremos que opines}

\begin{enumerate}
  \item \textbf{Errores de contenido.} Cualquier cosa que esté técnicamente mal,
        mal explicada o desactualizada. Es lo más importante.
  \item \textbf{Lo que no encaja con cómo trabajamos.} Buena parte de este
        material son \emph{propuestas}, no decisiones tomadas. En particular:
        \begin{itemize}
          \item la nomenclatura de piezas del apartado
                \ref{sec:herramientas_cad_datos:03} y la lista de códigos de
                subsistema;
          \item el sistema de coordenadas del apartado
                \ref{sec:herramientas_cad_datos:02};
          \item los pasos en \textcolor{fusalSec}{rojo discontinuo} de la figura
                \ref{fig:proceso_diseno}, que son añadidos al proceso de diseño y
                hay que decidir si se adoptan.
        \end{itemize}
        Si algo no te sirve, dilo ahora: cuesta mucho menos cambiarlo antes de
        que el equipo empiece a usarlo.
  \item \textbf{Lo que falta.} Temas que darías por imprescindibles en el tronco
        común y no están.
  \item \textbf{Si se entiende.} El público objetivo es alguien que acaba de
        entrar en el equipo. Si un apartado da por sabido algo que un rookie no
        sabe, señálalo.
\end{enumerate}

\section*{Sobre qué NO hace falta que opines}

\begin{itemize}
  \item \textbf{Los recuadros amarillos.} Son anotaciones de trabajo: cosas que ya
        sabemos que faltan. Están recogidas todas juntas al final del documento.
  \item \textbf{Los recuadros \textsc{caso fusal} que están sin rellenar.}
        Necesitan datos del equipo (puntuaciones, piezas rotas, masas) que
        todavía hay que recopilar.
  \item \textbf{Los capítulos \ref{ch:proceso_diseno} y
        \ref{ch:ensayo_adquisicion_datos}}, por lo dicho arriba. Del
        \ref{ch:proceso_diseno} sí interesa que revises las dos figuras.
  \item \textbf{La maquetación.} Portada, tipografía y colores están sin cerrar.
\end{itemize}

\section*{Cómo mandar los comentarios}

Indica siempre \textbf{número de capítulo y de apartado} (por ejemplo,
«\ref{ch:mecanica_aplicada}.4, lo del par de apriete») y qué cambiarías. Un
comentario sin localizar cuesta más de colocar que de escribir.

\vspace{0.8em}
\noindent
\begin{tabularx}{\textwidth}{@{}lX@{}}
\textbf{Enviar a:} & \pendiente{Rellenar: nombre y correo} \\[0.4em]
\textbf{Fecha límite:} & \pendiente{Rellenar} \\
\end{tabularx}

% !TeX root = ../../main.tex
% =====================================================================
%  Prólogo
% =====================================================================
\chapter*{Prólogo}
\addcontentsline{toc}{chapter}{Prólogo}
\markboth{Prólogo}{Prólogo}

En un equipo de Formula Student el problema nunca ha sido la falta de talento.
El problema es que cada uno o dos años se marcha quien sabía por qué se tomó
cada decisión, y el que llega detrás vuelve a empezar desde cero. A veces incluso
peor: repitiendo un error que ya se había cometido y resuelto tres temporadas
antes, sin que quedara escrito en ninguna parte.

Este libro existe para romper ese ciclo. No es un manual de teoría estricto --- para eso
están los libros que se citan en la bibliografía, y son mejores que este. Es el
registro de \emph{cómo diseñamos nuestro coche}: qué decidimos, con qué datos,
qué salió bien, qué salió mal y qué haríamos distinto. La teoría está aquí solo
en la medida necesaria para entender esas decisiones. Habrá conceptos en los que se pueda
profundizar más debido a que tiene un gran impacto en el diseñor y otros que se tratan de
manera superficial porque son o demasiado complejos para explicarlos aquí, o porque no son
tan relevantes para el diseño en el punto en el que está el equipo. En un futuro se espera
que se vaya mejorando el nivel de detalle y de introducir más resultados aclaratorios de
los ensayos que se realizan en el equipo. De momento y como primera versión, se priorizará
la claridad y utilidad de la información, y se irá mejorando con el tiempo.

De ahí las tres reglas con las que está escrito:

\begin{enumerate}
  \item \textbf{Todo capítulo termina en algo que se usa.} Una plantilla, una
        hoja de cálculo, un checklist, un apartado del informe de diseño. Si un
        capítulo no deja herramienta, sobra.
  \item \textbf{Ningún número sin origen.} Cada valor que aparezca debe poder
        rastrearse hasta un ensayo, un cálculo o una referencia. Los números sin
        procedencia son los que arruinan un design event.
  \item \textbf{Se escribe mientras se hace, no al final.} Un capítulo redactado
        seis meses después de la decisión ya ha perdido lo importante: las
        alternativas que se descartaron y por qué. (en la primera versión del libro, 
        esto no se ha cumplido del todo, pero se espera que en futuras versiones se 
        cumpla)
\end{enumerate}

Este libro nunca estará terminado, y eso es exactamente lo que se pretende. Cada
temporada debe salir una versión mejor que la anterior. Si encuentras algo mal,
desactualizado o simplemente incompleto, la respuesta correcta no es avisar a
alguien: es corregirlo y firmar el cambio.

\vspace{1.5em}
\noindent\hfill\textit{El equipo \FUSAL}

% !TeX root = ../../main.tex
% =====================================================================
%  Cómo usar este libro
% =====================================================================
\chapter*{Cómo usar este libro}
\addcontentsline{toc}{chapter}{Cómo usar este libro}
\markboth{Cómo usar este libro}{Cómo usar este libro}

\section*{Estructura}

El libro se divide en ocho partes. La primera es \textbf{tronco común}: la lee
todo el mundo, antes de tocar el CAD.
Las partes II a VII son los \textbf{itinerarios de cada subsistema}. La VIII es
\textbf{transversal} --- costes, fabricación, planificación --- y también es de
lectura obligatoria, porque el coste y la fabricabilidad son restricciones de
diseño, no un trámite del final de la temporada.

Las partes no son independientes. Como mínimo:

\begin{itemize}
  \item Dinámica vehicular (parte II) va \emph{antes} que suspensión (parte III):
        define los objetivos que la suspensión tiene que cumplir.
  \item La simulación de vuelta (capítulo \ref{ch:simulacion_vuelta}) es la
        herramienta con la que se justifican las decisiones de aerodinámica
        (parte V) y de relación de transmisión (parte VII).
  \item La refrigeración (parte VI) y la aerodinámica (parte V) comparten el
        mismo aire: los capítulos \ref{ch:lado_aire} y
        \ref{ch:elementos_aerodinamicos} se leen juntos.
\end{itemize}

\section*{Los recuadros}

Todos los capítulos usan los mismos seis recuadros. Debes poder hojear el libro
buscando solo un tipo de recuadro y sacar algo útil.

\begin{objetivos}
  \item Abre cada capítulo. Si al terminar no sabes hacer esto, vuelve atrás.
\end{objetivos}

\begin{clave}
La idea que hay que retener aunque se olvide todo lo demás del apartado.
\end{clave}

\begin{caso}[cómo se usa]
Lo que hicimos nosotros de verdad: la decisión, los datos de partida, el
resultado en pista y qué haríamos distinto. Es la parte del libro que no puedes
encontrar en ningún otro sitio, y por tanto la más importante.
\end{caso}

\begin{regla}[cómo se usa]
Requisito del reglamento que condiciona el diseño, con el artículo citado.
\textbf{Contrasta siempre con la versión vigente}: el reglamento cambia todos
los años y este libro no.
\end{regla}

\begin{ojo}
Los errores que se repiten cada temporada. Se escriben aquí precisamente para
dejar de repetirlos.
\end{ojo}

\begin{entregable}
Lo que produces al terminar el capítulo: plantilla, hoja de cálculo, checklist o
apartado del informe. Sin entregable, el capítulo no se da por leído.
\end{entregable}

\begin{juez}
Preguntas tipo del design event sobre el capítulo. Si no sabes contestarlas, no
has entendido el capítulo, independientemente de lo que creas.
\end{juez}

\section*{Cómo se imparte como curso}

Cada capítulo está pensado como una sesión de 90 minutos con este reparto:

\begin{itemize}
  \item \SI{20}{\minute} de teoría mínima --- solo lo necesario para entender el porqué.
  \item \SI{35}{\minute} sobre nuestro caso real (el recuadro \textsc{caso fusal}).
  \item \SI{30}{\minute} de ejercicio, produciendo el entregable.
  \item \SI{5}{\minute} de \emph{qué no sabemos todavía}: la lista de incógnitas
        abiertas que alimenta los TFG y los proyectos del año siguiente.
\end{itemize}

La sesión la imparten siempre \textbf{dos personas}: quien lo sabe y quien lo
tendrá que impartir el año que viene. Se graba siempre.

\section*{Responsables por capítulo}

\pendiente{Rellena esta tabla. Un capítulo sin responsable no se actualiza nunca.}

\begin{center}
\begin{tabular}{@{}llll@{}}
\toprule
\textbf{Parte} & \textbf{Responsable} & \textbf{Revisor} & \textbf{Últ. revisión} \\
\midrule
I --- Tronco común       & & & \\
II --- Dinámica vehicular & & & \\
III --- Suspensión        & & & \\
IV --- Chasis             & & & \\
V --- Aerodinámica        & & & \\
VI --- Refrigeración      & & & \\
VII --- Transmisión       & & & \\
VIII --- Transversal      & & & \\
\bottomrule
\end{tabular}
\end{center}

\section*{Qué leer además de esto}

Este libro es el registro de nuestras decisiones, no un sustituto de la teoría.
Estos son los libros a los que remite constantemente; conviene que el equipo
tenga al menos un ejemplar de cada uno.

\begin{description}[leftmargin=0pt,labelsep=0pt,style=nextline]
  \item[Dinámica vehicular] \textcite{milliken_rcvd_1995} es la referencia: densa,
    cara y con diferencia el mejor libro del campo. \textcite{gillespie_fundamentals_1992}
    es más asequible para empezar. Para el neumático,
    \textcite{pacejka_tirevehicle_2012}. \textcite{smith_tunetowin_1978} no tiene
    una sola ecuación y sigue explicando la puesta a punto mejor que casi nadie.
  \item[Suspensión y estructuras] \textcite{staniforth_suspension_2006} para la
    parte práctica y \textcite{smith_engineertowin_1984} para materiales y
    uniones. El dimensionado, con \textcite{budynas_shigley_2020}.
  \item[Aerodinámica] \textcite{katz_racecaraero_1995} está escrito justo para
    esto. Los fundamentos, en \textcite{anderson_aerodynamics_2017}; el CFD, en
    \textcite{versteeg_cfd_2007}.
  \item[Refrigeración] \textcite{incropera_heat_2011} para los fundamentos y
    \textcite{hesselgreaves_compact_2017} para intercambiadores compactos, que es
    lo que realmente montamos. \textcite{kays_london_1984} sigue siendo la fuente
    de los datos de superficie.
  \item[Formula Student] \textcite{seward_racecardesign_2014} recorre el coche
    entero con este nivel de detalle y es el mejor punto de partida para un
    rookie. Las notas técnicas de \textcite{optimumg_tips} son cortas y muy
    aplicables.
\end{description}

\section*{Cómo se edita}

\begin{itemize}
  \item Un fichero por capítulo, en \texttt{Capitulos/<parte>/}. No toques los
        ficheros de otros sin avisar.
  \item Mientras escribes un capítulo, descomenta \verb|\includeonly| en
        \texttt{main.tex} y deja solo el tuyo: compila en segundos.
  \item Las imágenes van en \texttt{Imagenes/<parte>/}, el código en
        \texttt{Codigo/}, los planos en PDF en \texttt{Planos/}.
  \item Símbolos: usa siempre las macros de \texttt{Estilo/macros.tex}
        ($\Cl$, $\Kus$, $\NTU$\ldots). Nunca los escribas a mano.
  \item Unidades: siempre con \texttt{siunitx} (\verb|\SI{85}{\celsius}|),
        nunca a pelo.
  \item Lo que falta se marca con \verb|\pendiente{}| o \verb|\falta{}| y sale
        listado al compilar en modo borrador. Nada de dejarlo «para acordarse».
\end{itemize}


\cleardoublepage
\tableofcontents

% !TeX root = ../../main.tex
% =====================================================================
%  Nomenclatura
%  Toda magnitud que aparezca en el libro debe estar en esta tabla.
%  Si añades un símbolo, añádelo también como macro en Estilo/macros.tex
% =====================================================================
\chapter*{Nomenclatura}
\addcontentsline{toc}{chapter}{Nomenclatura}
\markboth{Nomenclatura}{Nomenclatura}

Convenio de ejes: sistema SAE (\(x\) hacia delante, \(y\) hacia la derecha del
piloto, \(z\) hacia abajo), salvo indicación expresa en el capítulo
correspondiente. \pendiente{Confirma el convenio que usa el equipo y elimina esta nota.}

\section*{Dinámica vehicular}

\begin{longtable}{@{}p{2.2cm}p{8.5cm}p{2.5cm}@{}}
\toprule
\textbf{Símbolo} & \textbf{Magnitud} & \textbf{Unidad} \\
\midrule
\endfirsthead
\toprule
\textbf{Símbolo} & \textbf{Magnitud} & \textbf{Unidad} \\
\midrule
\endhead
$\Fz$      & Carga vertical sobre el neumático        & \si{\newton} \\
$\Fy$      & Fuerza lateral                            & \si{\newton} \\
$\Fx$      & Fuerza longitudinal                       & \si{\newton} \\
$\Mz$      & Par autoalineante                         & \si{\newton\meter} \\
$\slipa$   & Ángulo de deriva                          & \si{\degree} \\
$\slipr$   & Deslizamiento longitudinal                & --- \\
$\camber$  & Ángulo de caída                           & \si{\degree} \\
$\muy$     & Coeficiente de agarre lateral             & --- \\
$\ay$      & Aceleración lateral                       & \si{\meter\per\second\squared} \\
$\ax$      & Aceleración longitudinal                  & \si{\meter\per\second\squared} \\
$\hcg$     & Altura del centro de gravedad             & \si{\milli\meter} \\
$\batalla$ & Batalla                                   & \si{\milli\meter} \\
$\via$     & Ancho de vía                              & \si{\milli\meter} \\
$\Kroll$   & Rigidez a balanceo                        & \si{\newton\meter\per\degree} \\
$\LLTD$    & Reparto de transferencia de carga lateral & \si{\percent} \\
$\Kus$     & Gradiente de subviraje                    & \si{\degree\per\gacc} \\
\bottomrule
\end{longtable}

\section*{Suspensión}

\begin{longtable}{@{}p{2.2cm}p{8.5cm}p{2.5cm}@{}}
\toprule
\textbf{Símbolo} & \textbf{Magnitud} & \textbf{Unidad} \\
\midrule
\endfirsthead
\toprule
\textbf{Símbolo} & \textbf{Magnitud} & \textbf{Unidad} \\
\midrule
\endhead
$\MR$ & Relación de instalación (motion ratio) & --- \\
$\Kw$ & Rigidez medida en rueda                 & \si{\newton\per\milli\meter} \\
$\Ks$ & Rigidez del muelle                      & \si{\newton\per\milli\meter} \\
$\fr$ & Frecuencia de suspensión                & \si{\hertz} \\
\bottomrule
\end{longtable}

\section*{Aerodinámica}

\begin{longtable}{@{}p{2.2cm}p{8.5cm}p{2.5cm}@{}}
\toprule
\textbf{Símbolo} & \textbf{Magnitud} & \textbf{Unidad} \\
\midrule
\endfirsthead
\toprule
\textbf{Símbolo} & \textbf{Magnitud} & \textbf{Unidad} \\
\midrule
\endhead
$\Cl$    & Coeficiente de sustentación (negativo = carga)     & --- \\
$\Cd$    & Coeficiente de resistencia                          & --- \\
$\Cp$    & Coeficiente de presión                              & --- \\
$\ClA$   & Producto coeficiente-área de carga                  & \si{\meter\squared} \\
$\CdA$   & Producto coeficiente-área de resistencia            & \si{\meter\squared} \\
$\qdin$  & Presión dinámica                                    & \si{\pascal} \\
$\Vinf$  & Velocidad de la corriente libre                     & \si{\meter\per\second} \\
$\Rey$   & Número de Reynolds                                  & --- \\
$\CoP$   & Centro de presiones (\% de batalla desde el eje delantero) & \si{\percent} \\
\bottomrule
\end{longtable}

\section*{Transferencia de calor}

\begin{longtable}{@{}p{2.2cm}p{8.5cm}p{2.5cm}@{}}
\toprule
\textbf{Símbolo} & \textbf{Magnitud} & \textbf{Unidad} \\
\midrule
\endfirsthead
\toprule
\textbf{Símbolo} & \textbf{Magnitud} & \textbf{Unidad} \\
\midrule
\endhead
$\Qdot$  & Potencia térmica a disipar                & \si{\watt} \\
$\mdot$  & Gasto másico                              & \si{\kilogram\per\second} \\
$\UA$    & Conductancia global del intercambiador    & \si{\watt\per\kelvin} \\
$\NTU$   & Número de unidades de transferencia       & --- \\
$\efec$  & Eficacia del intercambiador               & --- \\
$\LMTD$  & Diferencia de temperatura media logarítmica & \si{\kelvin} \\
$\jcol$  & Factor de Colburn                          & --- \\
$\ffan$  & Factor de fricción                         & --- \\
$\Nus$   & Número de Nusselt                          & --- \\
$\Pr$    & Número de Prandtl                          & --- \\
$\Dh$    & Diámetro hidráulico                        & \si{\milli\meter} \\
\bottomrule
\end{longtable}

\section*{Transmisión}

\begin{longtable}{@{}p{2.2cm}p{8.5cm}p{2.5cm}@{}}
\toprule
\textbf{Símbolo} & \textbf{Magnitud} & \textbf{Unidad} \\
\midrule
\endfirsthead
\toprule
\textbf{Símbolo} & \textbf{Magnitud} & \textbf{Unidad} \\
\midrule
\endhead
$\itot$  & Relación de transmisión total & --- \\
$\Tmot$  & Par motor                     & \si{\newton\meter} \\
$\rdin$  & Radio dinámico del neumático  & \si{\milli\meter} \\
\bottomrule
\end{longtable}

\section*{Acrónimos}

\begin{longtable}{@{}p{2.6cm}p{10.5cm}@{}}
\toprule
\textbf{Acrónimo} & \textbf{Significado} \\
\midrule
\endfirsthead
\toprule
\textbf{Acrónimo} & \textbf{Significado} \\
\midrule
\endhead
AMFE   & Análisis Modal de Fallos y Efectos (FMEA) \\
BOM    & \textit{Bill of Materials}, lista de materiales \\
CFD    & \textit{Computational Fluid Dynamics} \\
CdG    & Centro de gravedad \\
CoP    & \textit{Centre of Pressure}, centro de presiones \\
DFM    & \textit{Design for Manufacturing} \\
FEA    & \textit{Finite Element Analysis} \\
GGV    & Envolvente de aceleraciones frente a velocidad \\
IAD    & \textit{Impact Attenuator Data}, informe del atenuador de impactos \\
KPI    & \textit{Kingpin Inclination}, inclinación del eje de pivote \\
LLTD   & \textit{Lateral Load Transfer Distribution} \\
LMTD   & \textit{Log Mean Temperature Difference} \\
MMM    & \textit{Milliken Moment Method} \\
NTU    & \textit{Number of Transfer Units} \\
SES    & \textit{Structural Equivalency Spreadsheet} \\
TTC    & \textit{Tire Test Consortium} \\
YMD    & \textit{Yaw Moment Diagram} \\
\bottomrule
\end{longtable}


% =====================================================================
\mainmatter
% =====================================================================

% ---------------------------------------------------------------------
\part{Tronco común}
% ---------------------------------------------------------------------
% !TeX root = ../../main.tex
% =====================================================================
%  Formula Student: el juego y sus reglas
%  Responsable:
%  Última revisión:
% =====================================================================
\chapter{Formula Student: el juego y sus reglas}
\label{ch:formula_student}

\begin{objetivos}
  \item Enumerar los eventos de una competición y cuántos puntos vale cada uno.
  \item Explicar por qué la fiabilidad da más puntos que las prestaciones.
  \item Localizar en el reglamento el artículo que afecta a tu subsistema y saber
        qué hacer cuando es ambiguo.
  \item Reconocer los documentos de temporada y las consecuencias de entregarlos tarde.
  \item Ordenar por prioridad el trabajo del equipo en función de dónde están los puntos.
\end{objetivos}

\fichasesion{90 minutos}{Ninguno. Es la primera sesión del curso}%
{Reglamento vigente descargado y calendario de la temporada}

Formula Student no es una carrera. Es una competición de \emph{ingeniería} en la
que hay, además, unas pruebas en pista. Esta distinción no es un matiz retórico:
determina cómo se reparten los puntos, y por tanto cómo debe repartir el equipo
su tiempo. Un tercio de la puntuación se decide hablando con jueces, sin
encender el coche.

Este capítulo describe cómo se reparten los puntos y qué documentos hay que
entregar durante la temporada. El resto del libro se apoya en ello: las
decisiones de diseño de los capítulos siguientes se justifican, en última
instancia, contra este reparto.

% ---------------------------------------------------------------------
\section{La competición y sus eventos}
\label{sec:formula_student:01}
\index{eventos}

Una competición se divide en \textbf{eventos estáticos} y \textbf{eventos
dinámicos}, con una barrera entre medias: la \textbf{revisión técnica}
(\emph{scrutineering}).

\subsection{Eventos estáticos}

Se disputan con el coche parado, delante de un panel de jueces que son
ingenieros del sector.

\begin{description}[leftmargin=0pt,labelsep=0pt,style=nextline]
  \item[\textit{Business Plan Presentation}] El equipo presenta un plan de
    negocio para el vehículo como si fuera un producto real. No se juzga la
    ingeniería, sino la coherencia del modelo de negocio y la calidad de la
    presentación.
  \item[\textit{Cost and Manufacturing}] Se entrega un informe con el coste
    completo del coche calculado según unas tablas comunes a todos los equipos, y
    se defiende ante los jueces. Se evalúa la precisión del informe, el
    conocimiento real de los procesos de fabricación y, en el \textit{Real Case
    Scenario}, la capacidad de razonar sobre coste bajo un supuesto que dan en el
    momento. Se trata en detalle en el capítulo \ref{ch:evento_costes}.
  \item[\textit{Engineering Design}] El evento más importante de los tres y el
    que más puntos reparte. Los jueces preguntan \emph{por qué} está diseñado así
    cada subsistema. No buscan el coche más sofisticado: buscan decisiones
    justificadas con datos y coherentes con los objetivos que el propio equipo se
    ha fijado. Se prepara con el capítulo \ref{ch:design_event}.
\end{description}

\subsection{Revisión técnica}
\index{revisión técnica}\index{scrutineering}

No reparte puntos. Es una \emph{puerta}: hasta que el coche no la pasa, no puede
salir a pista. Incluye, según el tipo de vehículo y la competición, la inspección
mecánica y eléctrica, la prueba de vuelco, el ensayo de ruido o de estanqueidad,
la prueba de salida del piloto (\emph{egress}) y el ensayo de frenada.

\begin{clave}
Superar la revisión técnica vale, en la práctica, los \textbf{675 puntos
dinámicos}: sin la pegatina, todos ellos son cero. Es la actividad con mayor
rentabilidad por hora invertida de toda la temporada.
\end{clave}

\subsection{Eventos dinámicos}

\begin{description}[leftmargin=0pt,labelsep=0pt,style=nextline]
  \item[\textit{Acceleration}] Recta de \SI{75}{\meter} desde parado. Mide
    tracción y relación de transmisión; el aerodinámico aquí solo estorba.
  \item[\textit{Skidpad}] Un ocho de radio constante. Mide agarre lateral
    estacionario y equilibrio del coche.
  \item[\textit{Autocross}] Una vuelta rápida a un circuito estrecho y revirado.
    Además de puntuar, fija el orden de salida de la resistencia.
  \item[\textit{Endurance}] \SI{22}{\kilo\meter} con un cambio de piloto a
    mitad. Es el evento decisivo: mide fiabilidad y, de paso, ritmo sostenido.
  \item[\textit{Efficiency}] No se corre: se calcula con el consumo y el tiempo
    de la resistencia. \emph{Solo puntúa si se termina la resistencia.}
\end{description}

% ---------------------------------------------------------------------
\section{El reparto de puntos}
\label{sec:formula_student:02}
\index{puntuación}

\begin{table}[H]
\centering
\caption{Reparto habitual de la puntuación. Verificar siempre contra el
reglamento vigente: los valores cambian entre categorías y entre temporadas.}
\label{tab:reparto_puntos}
\begin{tabular}{@{}llS[table-format=3.0]@{}}
\toprule
\textbf{Evento} & \textbf{Tipo} & {\textbf{Puntos}} \\
\midrule
Business Plan Presentation & Estático & 75 \\
Cost and Manufacturing     & Estático & 100 \\
Engineering Design         & Estático & 150 \\
\midrule
\multicolumn{2}{@{}l}{\textbf{Subtotal estáticos}} & \bfseries 325 \\
\midrule
Acceleration & Dinámico & 75 \\
Skidpad      & Dinámico & 75 \\
Autocross    & Dinámico & 100 \\
Endurance    & Dinámico & 325 \\
Efficiency   & Dinámico & 100 \\
\midrule
\multicolumn{2}{@{}l}{\textbf{Subtotal dinámicos}} & \bfseries 675 \\
\midrule
\multicolumn{2}{@{}l}{\textbf{Total}} & \bfseries 1000 \\
\bottomrule
\end{tabular}
\end{table}

Conviene mirar esta tabla con atención, porque contiene casi toda la estrategia
del equipo:

\begin{itemize}
  \item \textbf{Los estáticos son 325 puntos que no dependen de que el coche
        funcione.} Se pueden preparar con el coche a medio montar, y de hecho es
        lo que hacen los equipos que puntúan bien en ellos.
  \item \textbf{\textit{Endurance} y \textit{Efficiency} suman 425 puntos}, más
        que todos los estáticos juntos. Y están acoplados: sin terminar la
        resistencia, ambos son cero.
  \item \textbf{\textit{Acceleration} y \textit{Skidpad} suman 150 puntos.} Son
        los eventos donde más se nota el rendimiento puro y, sin embargo, los que
        menos pesan. Es un error clásico diseñar el coche para ellos.
\end{itemize}

Además, la puntuación de cada evento dinámico no es lineal: se reparte por
comparación con el mejor equipo y con un tiempo máximo, de modo que
\textbf{terminar despacio puntúa mucho más que no terminar}. La diferencia entre
el primero y el último clasificado \emph{de los que terminan} es siempre menor
que la diferencia entre terminar y no terminar.

% ---------------------------------------------------------------------
\section{Cómo se lee el reglamento}
\label{sec:formula_student:03}
\index{reglamento}

El reglamento no es un anexo administrativo: es el primer documento de
requisitos del proyecto, y llega escrito antes de que empieces a diseñar. Buena
parte de las decisiones de arquitectura del coche están ya tomadas en él.

\subsection{Estructura}

Las normas están agrupadas por secciones identificadas con letras, y cada
artículo se cita con la letra de su sección y una numeración jerárquica (por
ejemplo, \texttt{T 2.1.1}). Las secciones habituales son:

\begin{center}
\begin{tabular}{@{}ll@{}}
\toprule
\textbf{Sección} & \textbf{Contenido} \\
\midrule
A  & Normas administrativas: inscripción, plazos, documentos, penalizaciones \\
T  & Normas técnicas generales, aplicables a todos los vehículos \\
C  & Específicas de vehículos de combustión \\
EV & Específicas de vehículos eléctricos \\
DV & Específicas de vehículos autónomos \\
IN & Inspecciones y revisión técnica \\
S  & Eventos estáticos \\
D  & Eventos dinámicos \\
\bottomrule
\end{tabular}
\end{center}

\begin{regla}[cómo citar]
En este libro, y en cualquier documento del equipo, un requisito normativo se
cita \emph{siempre} con su artículo y la versión del reglamento:
«\textit{la distancia mínima entre el casco y la línea que une los arcos es de
\SI{50}{\milli\meter} (T 1.4.2, FS Rules 2025 v1.1)}». Sin artículo no es un
requisito: es una opinión.
\end{regla}

\subsection{Método de trabajo}

\begin{enumerate}
  \item \textbf{Léelo entero una vez, al principio de la temporada}, incluidas
        las secciones que no son de tu subsistema. Muchos requisitos afectan a
        varios a la vez.
  \item \textbf{Extrae los artículos de tu subsistema a una tabla de
        trazabilidad} con cuatro columnas: artículo, qué exige, cómo lo
        cumplimos, dónde se demuestra. Esa tabla es simultáneamente la
        preparación de la revisión técnica y media memoria de diseño.
  \item \textbf{Distingue lo obligatorio de lo recomendado.} El reglamento usa un
        lenguaje deliberadamente preciso; una frase con «\textit{must}» o
        «\textit{shall}» no admite interpretación.
  \item \textbf{Si algo es ambiguo, pregunta por el cauce oficial.} Las
        competiciones tienen un sistema de consultas cuyas respuestas son
        vinculantes. Una respuesta por escrito vale en la revisión técnica; la
        interpretación de un compañero, no.
  \item \textbf{No copies números de la temporada anterior.} El reglamento cambia
        todos los años, y las cotas, espesores y materiales admisibles están
        entre lo que más cambia.
\end{enumerate}

\begin{ojo}
\begin{itemize}
  \item Diseñar primero y buscar el artículo después, para «comprobar que
        cumple». El orden correcto es el inverso.
  \item Leer solo la sección propia. Los requisitos de accesibilidad,
        protecciones y separaciones son transversales y se incumplen en las
        \emph{fronteras} entre subsistemas.
  \item Dar por bueno el diseño del coche anterior. Que pasara la revisión hace
        dos temporadas no significa que la pase con el reglamento de esta.
\end{itemize}
\end{ojo}

% ---------------------------------------------------------------------
\section{Documentos y plazos de la temporada}
\label{sec:formula_student:04}
\index{documentos de temporada}

Una parte importante de la competición se juega meses antes de llegar al
circuito, entregando documentos. Casi todos tienen penalización por retraso, y
algunos son excluyentes: sin ellos no se compite.

\begin{table}[H]
\centering
\caption{Documentos habituales de temporada. Los plazos concretos y la lista
exacta dependen de la competición: rellenar al inicio de cada temporada.}
\label{tab:documentos_temporada}
\begin{tabularx}{\textwidth}{@{}lXll@{}}
\toprule
\textbf{Documento} & \textbf{Qué es} & \textbf{Responsable} & \textbf{Fecha} \\
\midrule
Inscripción y cuestionario & Prueba de conocimiento del reglamento & & \\
SES & Justificación de equivalencia estructural del chasis & & \\
IAD & Diseño y ensayo del atenuador de impactos & & \\
ESF & Formulario del sistema eléctrico (vehículos eléctricos) & & \\
Cost Report y BOM & Coste completo del vehículo & & \\
Engineering Design Report & Memoria de diseño & & \\
Design Spec Sheet & Ficha resumen de características & & \\
Business Plan & Documentación del evento de negocio & & \\
Vídeo de estado del vehículo & Prueba de que el coche rueda & & \\
\bottomrule
\end{tabularx}
\end{table}

\pendiente{Rellenar la tabla con los documentos, responsables y fechas reales de
la temporada en curso, y colgarla donde se vea todos los días.}

\begin{clave}
Los plazos de documentación son la única parte de la competición donde se pierden
puntos \emph{sin haber hecho nada mal técnicamente}. Cuestan cero horas de taller
y solo exigen organización. Un equipo que pierde puntos aquí está regalando la
parte más barata de la competición.
\end{clave}

% ---------------------------------------------------------------------
\section{De dónde salen realmente los puntos}
\label{sec:formula_student:05}

Si se cruza el reparto de puntos con lo que ocurre de verdad en una competición,
sale una lista de prioridades bastante distinta de la intuitiva. Un porcentaje
muy alto de los equipos inscritos no termina la resistencia, y una parte
importante ni siquiera llega a pasar la revisión técnica a tiempo de disputar
todos los eventos.

De ahí el orden de prioridad con el que debería trabajar cualquier equipo, y muy
especialmente uno joven:

\begin{enumerate}
  \item \textbf{Pasar la revisión técnica.} Diseñar desde el primer día con la
        lista de comprobación de la inspección delante, y hacer una inspección
        interna simulada semanas antes de viajar.
  \item \textbf{Terminar la resistencia.} Kilómetros de rodaje antes de la
        competición, componentes con margen, repuestos y procedimientos de
        montaje. La fiabilidad no se añade al final: se diseña.
  \item \textbf{Preparar bien los estáticos.} 325 puntos que solo cuestan
        organización, documentación y ensayo de la defensa.
  \item \textbf{Y después, ir rápido.} Cuando lo anterior está resuelto, las
        prestaciones deciden posiciones. Antes, no deciden nada.
\end{enumerate}

\begin{clave}
El coche no tiene que ser rápido: tiene que \textbf{terminar}. Un coche lento que
completa los cinco eventos dinámicos y defiende bien sus estáticos puntúa por
encima de un coche brillante que rompe en la resistencia.
\end{clave}

Esta jerarquía tiene una consecuencia incómoda para un equipo de ingenieros: casi
siempre, la decisión correcta es \emph{la más simple que cumple el objetivo}. La
sofisticación solo se justifica cuando el equipo ya es capaz de fabricarla,
validarla y repararla en el paddock. Los jueces del evento de diseño valoran
exactamente eso, y penalizan la complejidad que el equipo no puede sostener.

% ---------------------------------------------------------------------
%  Cierre del capítulo
% ---------------------------------------------------------------------

\begin{caso}
\redactar{Rellenar con nuestro historial: puntuación obtenida por evento en las
últimas participaciones, en qué eventos no pudimos puntuar y por qué, y cuánto
tiempo tardamos en pasar la revisión técnica. Es el punto de partida honesto
para fijar los objetivos de esta temporada.}
\end{caso}

\begin{ojo}
\begin{itemize}
  \item Medir el éxito del equipo por el tiempo en \textit{acceleration}: 75 de
        1000 puntos.
  \item Empezar a preparar los estáticos cuando el coche ya rueda. Para entonces
        no queda tiempo, y son 325 puntos.
  \item Tratar la revisión técnica como un trámite del final en lugar de como un
        requisito de diseño.
  \item Fijar los objetivos de la temporada sin mirar la puntuación de la
        anterior.
\end{itemize}
\end{ojo}

\begin{entregable}
\begin{enumerate}
  \item La tabla \ref{tab:documentos_temporada} rellena con los plazos reales de
        esta temporada y un responsable por documento.
  \item Una tabla de trazabilidad de reglamento para tu subsistema: artículo,
        qué exige, cómo lo cumplimos, dónde se demuestra.
  \item Los objetivos de puntuación del equipo para esta temporada, evento a
        evento, con su justificación.
\end{enumerate}
\end{entregable}

\begin{juez}
\begin{enumerate}
  \item ¿Qué objetivos de puntuación se fijó el equipo esta temporada y por qué esos?
  \item ¿Qué evento habéis priorizado en el diseño del coche, y qué habéis
        sacrificado a cambio?
  \item ¿Qué artículo del reglamento ha condicionado más el diseño de tu
        subsistema? Cítalo.
  \item ¿Cuántos kilómetros ha rodado el coche antes de venir aquí?
  \item ¿Qué os impidió puntuar en algún evento la temporada pasada y qué habéis
        cambiado para que no vuelva a pasar?
\end{enumerate}
\end{juez}

% !TeX root = ../../main.tex
% =====================================================================
%  El coche como sistema
%  Responsable:
%  Última revisión:
%  NOTA: capítulo deliberadamente breve. Describe el método al que el equipo
%  quiere llegar, no el que ha seguido esta temporada. Ver el recuadro CASO.
% =====================================================================
\chapter{El coche como sistema}
\label{ch:coche_como_sistema}

\begin{objetivos}
  \item Distinguir un objetivo de vehículo de un deseo.
  \item Explicar cómo un objetivo global se convierte en un requisito concreto de
        una pieza.
  \item Entender qué es un presupuesto de masa y para qué sirve un modelo de
        centro de gravedad.
  \item Explicar por qué la simulación de vuelta es el único árbitro razonable
        entre subsistemas que compiten por lo mismo.
\end{objetivos}

\fichasesion{60 minutos}{Capítulo \ref{ch:formula_student}}%
{Hoja de cálculo de masas del coche actual, si existe}

Este capítulo es corto a propósito, y conviene decir por qué desde el principio:
describe la forma de trabajar a la que el equipo quiere llegar, no la que ha
seguido hasta ahora. Se incluye en el tronco común porque es el marco que da
sentido a los capítulos de dinámica vehicular, aerodinámica y transmisión, y
porque los jueces del evento de diseño preguntan por él antes que por cualquier
detalle de una pieza.

La idea de fondo es simple. Un coche de competición es un sistema en el que casi
ningún subsistema puede mejorar sin empeorar otro: el aerodinámico añade carga y
también masa y resistencia; el radiador necesita aire que la carrocería quería
para otra cosa; una relación de transmisión corta gana en aceleración y pierde en
velocidad punta. Si cada subsistema optimiza lo suyo por separado, el resultado no
es un coche optimizado, sino un coche incoherente. Hace falta un criterio común
--- unos objetivos de vehículo --- y una herramienta que traduzca cualquier
decisión local a ese criterio.

% ---------------------------------------------------------------------
\section{Objetivos de vehículo}
\label{sec:coche_como_sistema:01}
\index{objetivos de vehículo}

Un objetivo de vehículo es un número, con fecha y con responsable. «El coche tiene
que ser ligero» no es un objetivo; «masa en orden de marcha sin piloto por debajo
de \SI{230}{\kilogram}, verificada en báscula antes del 15 de mayo, responsable
el jefe de chasis» sí lo es.

Los objetivos típicos de un coche de Formula Student son media docena:

\begin{center}
\begin{tabularx}{\textwidth}{@{}lXl@{}}
\toprule
\textbf{Objetivo} & \textbf{Por qué se fija} & \textbf{Unidad} \\
\midrule
Masa total & Afecta a todo: aceleración, frenada, curva, consumo & \si{\kilogram} \\
Reparto de masas & Condiciona el equilibrio y la tracción & \si{\percent} \\
Altura del centro de gravedad & Gobierna la transferencia de carga & \si{\milli\meter} \\
Rigidez torsional del chasis & Determina si la suspensión hace lo que se diseñó & \si{\newton\meter\per\degree} \\
Carga aerodinámica y resistencia & Prestaciones en curva frente a recta & \si{\meter\squared} \\
Tiempo objetivo por vuelta & Integra todo lo anterior en un solo número & \si{\second} \\
\bottomrule
\end{tabularx}
\end{center}

Lo importante no es acertar el valor a la primera --- casi nunca se acierta ---
sino que exista, que esté escrito y que cuando se incumpla \emph{se sepa}. Si un
objetivo se revisa a la baja, la revisión se discute y se anota.

\begin{clave}
El valor de un objetivo de vehículo no está en el número, sino en que obliga a
tener la discusión \emph{antes} de fabricar. Cuando dos subsistemas piden lo
mismo, el objetivo es lo que permite decidir sin que gane el que más levanta la voz.
\end{clave}

% ---------------------------------------------------------------------
\section{Cascada de requisitos}
\label{sec:coche_como_sistema:02}
\index{cascada de requisitos}

Un objetivo de vehículo no se puede diseñar directamente: hay que descomponerlo
hasta llegar a algo que una persona pueda dimensionar. Ese descenso es la cascada
de requisitos, y tiene tres niveles:

\begin{enumerate}
  \item \textbf{Objetivo de vehículo.} Masa total por debajo de
        \SI{230}{\kilogram}.
  \item \textbf{Requisito de subsistema.} Suspensión completa (cuatro esquinas)
        por debajo de \SI{28}{\kilogram}.
  \item \textbf{Requisito de componente.} Trapecio delantero superior por debajo
        de \SI{380}{\gram} cumpliendo los casos de carga del capítulo
        \ref{ch:casos_carga}.
\end{enumerate}

La cascada funciona en los dos sentidos. Hacia abajo reparte el presupuesto; hacia
arriba avisa. Si el trapecio se va a \SI{450}{\gram} y no hay forma de bajarlo,
alguien tiene que decidir conscientemente si se acepta, si se compensa en otro
sitio o si se cambia el objetivo. Lo que no puede pasar es que nadie se entere
hasta la báscula.

% ---------------------------------------------------------------------
\section{Presupuesto de masa y centro de gravedad}
\label{sec:coche_como_sistema:03}
\index{presupuesto de masa}\index{centro de gravedad}

El presupuesto de masa es una hoja de cálculo con una fila por pieza y tres
columnas de masa: \textbf{estimada} (al fijar el objetivo), \textbf{de CAD}
(cuando la pieza está modelada, con el material asignado de verdad) y
\textbf{pesada} (cuando existe). La diferencia entre la primera y la última
columna indica cómo de fiables son nuestras estimaciones, y sirve para ajustarlas
la temporada siguiente.

Si además se anota la posición del centro de gravedad de cada pieza, la misma
hoja da el centro de gravedad del coche completo:

\begin{equation}
  x_{\CdG} = \frac{\sum_i m_i\, x_i}{\sum_i m_i}
  \label{eq:cdg}
\end{equation}

y lo mismo para $y$ y $z$. No hace falta nada más sofisticado. Con eso se puede
responder a preguntas del tipo «¿cuánto sube el centro de gravedad si monto el
radiador aquí en vez de allí?» en el momento en que se plantean, que es cuando
sirve de algo.

La altura del centro de gravedad merece atención especial porque gobierna la
transferencia de carga (capítulo \ref{ch:regimen_estacionario}) y, con ella, el
agarre disponible. Suele recibir mucha menos atención que la masa total, y en un
coche de Formula Student unos pocos centímetros de altura pesan tanto como varios
kilogramos.

\begin{ojo}
\begin{itemize}
  \item Actualizar el presupuesto de masa solo un par de veces al año: deja de
        servir para decidir. Se actualiza cada vez que se libera una pieza.
  \item Estimar masas «por experiencia» sin apuntar de dónde sale la estimación:
        cuando falla, no se puede aprender nada.
  \item Olvidar en la hoja lo que no es una pieza: cableado, líquidos, tornillería,
        pintura, cinta. Suele ser bastante más de lo que nadie espera.
  \item Medir el centro de gravedad solo en el plano y no en altura, que es lo
        que realmente condiciona el comportamiento.
\end{itemize}
\end{ojo}

% ---------------------------------------------------------------------
\section{La simulación de vuelta como árbitro}
\label{sec:coche_como_sistema:04}
\index{simulación de vuelta}

Supongamos la discusión clásica: el paquete aerodinámico añade
\SI{0.9}{\meter\squared} de $\ClA$ y \SI{12}{\kilogram} de masa. ¿Compensa? No
hay forma de contestar razonando: hay que calcularlo, porque una parte del efecto
es favorable y otra no, y el resultado depende del trazado, de las velocidades y
del agarre disponible.

Eso es exactamente lo que hace una simulación de vuelta: convierte un cambio de
diseño en segundos por vuelta y, a través de la tabla de puntos del capítulo
\ref{ch:formula_student}, en puntos. Su valor real no está en predecir el tiempo
absoluto --- eso casi nunca lo acierta --- sino en dar \textbf{sensibilidades}:

\begin{center}
\begin{tabular}{@{}ll@{}}
\toprule
\textbf{Sensibilidad} & \textbf{Unidades} \\
\midrule
Segundos por vuelta por kilogramo & \si{\second\per\kilogram} \\
Segundos por vuelta por unidad de $\ClA$ & \si{\second\per\meter\squared} \\
Segundos por vuelta por milímetro de altura de $\CdG$ & \si{\second\per\milli\meter} \\
Segundos por vuelta por punto de coeficiente de agarre & \si{\second} \\
\bottomrule
\end{tabular}
\end{center}

Con esa tabla, la discusión anterior se resuelve con un número. Y, de cara al
evento de diseño, la respuesta a «¿por qué lleváis alerón?» deja de ser una
opinión.

Estas sensibilidades son la entrada de los capítulos
\ref{ch:merece_la_pena_aero} (aerodinámica) y \ref{ch:relacion_transmision}
(relación de transmisión). Sin ellas, esos dos capítulos se quedan en criterio
cualitativo.

\begin{clave}
La simulación de vuelta no sirve para saber cuánto se tarda: sirve para saber
\textbf{cuánto cambia el tiempo} cuando cambias algo. Un modelo mediocre bien
entendido da sensibilidades útiles; un modelo excelente sin validar da un número
absoluto en el que nadie debería confiar.
\end{clave}

% ---------------------------------------------------------------------
\section{Interfaces entre subsistemas}
\label{sec:coche_como_sistema:05}
\index{interfaces}

Los fallos de integración no ocurren dentro de los subsistemas, sino en las
fronteras, y aparecen tarde: en el montaje. Se previenen con una tabla de
interfaces que se acuerda al principio de la temporada y en la que cada línea
tiene un dueño único.

\begin{center}
\begin{tabularx}{\textwidth}{@{}llX@{}}
\toprule
\textbf{Interfaz} & \textbf{Dueño} & \textbf{Qué hay que acordar} \\
\midrule
Chasis--suspensión & Chasis & Posición y tolerancia de los anclajes \\
Suspensión--freno--rueda & Suspensión & Espacio libre en la llanta, cotas de buje \\
Refrigeración--aerodinámica & Aerodinámica & Caudal y superficie de entrada de aire \\
Chasis--transmisión & Transmisión & Anclajes, alineación y resguardos \\
Chasis--ergonomía & Ergonomía & Volumen de cockpit, plantillas de reglamento \\
Electrónica--todos & Electrónica & Recorrido de mazos, sensores, masas \\
\bottomrule
\end{tabularx}
\end{center}

La regla práctica es que \textbf{ninguna interfaz puede tener dos dueños}. Si dos
subsistemas comparten una cota, uno la fija y el otro la consume; y la cota vive
en el esqueleto del CAD (capítulo \ref{ch:herramientas_cad_datos}), no en la
cabeza de nadie.

% ---------------------------------------------------------------------
%  Cierre del capítulo
% ---------------------------------------------------------------------

\begin{caso}[dónde estamos de verdad]
Esta temporada el equipo \textbf{no ha trabajado con objetivos de vehículo
formales}: no se fijó un objetivo de masa, ni de altura de centro de gravedad, ni
un tiempo por vuelta de referencia, y las decisiones de cada subsistema se han
tomado con criterio local. Tampoco se ha usado una simulación de vuelta para
arbitrar entre subsistemas.

Conviene escribirlo tal cual, por dos razones. La primera es que este libro solo
sirve si es honesto. La segunda es que \emph{es exactamente la pregunta que hacen
los jueces de diseño}: cómo se decidió el reparto de prioridades del coche. Tener
clara la respuesta ---incluso cuando la respuesta es «este año no lo hicimos
así, y estas son las consecuencias que hemos visto»--- puntúa más que improvisar
una justificación.

El camino para la próxima temporada es el que describe este capítulo, en este orden:

\begin{enumerate}
  \item Montar el presupuesto de masa con lo que ya existe, aunque sea a
        posteriori: pesar el coche actual pieza a pieza. Está comprometido y
        pendiente de hacer; es la base de todo lo demás.
  \item Medir el centro de gravedad del coche actual, en las tres direcciones.
  \item Levantar una simulación de vuelta sencilla, aunque sea cuasi-estacionaria
        (capítulo \ref{ch:simulacion_vuelta}), y sacar las cuatro sensibilidades
        de la tabla anterior.
  \item Con esos datos ya sí: fijar los objetivos de la temporada siguiente.
\end{enumerate}

\redactar{Actualizar este recuadro cuando se completen los pasos anteriores,
dejando constancia de las cifras reales obtenidas.}
\end{caso}

\begin{ojo}
\begin{itemize}
  \item Fijar objetivos de vehículo copiando los de otro equipo. Un objetivo que
        el equipo no sabe de dónde sale no se defiende ante un juez ni se cumple.
  \item Poner objetivos tan holgados que se cumplen solos: no informan de nada.
  \item Usar la simulación de vuelta para decidir y no validarla nunca contra
        datos de pista.
  \item Descubrir las interfaces en el montaje.
\end{itemize}
\end{ojo}

\begin{entregable}
\begin{enumerate}
  \item La hoja de presupuesto de masa del coche actual, con las columnas de masa
        estimada, de CAD y pesada.
  \item La tabla de interfaces con un dueño por línea, acordada y firmada por los
        responsables de subsistema.
\end{enumerate}
\end{entregable}

\begin{juez}
\begin{enumerate}
  \item ¿Qué objetivos de vehículo os fijasteis y de dónde salieron esos valores?
  \item ¿Cuánto pesa el coche y cuánto pensabais que iba a pesar al empezar a
        diseñarlo?
  \item ¿Dónde está el centro de gravedad y cómo lo habéis medido?
  \item ¿Cuántos segundos por vuelta os cuesta cada kilogramo?
  \item Cuando aerodinámica y refrigeración han pedido el mismo espacio, ¿quién
        ha decidido y con qué criterio?
\end{enumerate}
\end{juez}

% !TeX root = ../../main.tex
% =====================================================================
%  Proceso de diseño y toma de decisiones
%  Responsable: 
%  Última revisión: 
% =====================================================================
\chapter{Proceso de diseño y toma de decisiones}
\label{ch:proceso_diseno}

\begin{objetivos}
  \item \redactar{Objetivo de aprendizaje 1 --- concreto y comprobable.}
  \item \redactar{Objetivo 2.}
  \item \redactar{Objetivo 3.}
\end{objetivos}

\fichasesion{90 min}{\redactar{capítulos previos}}{\redactar{material}}

\section{Fases del ciclo de diseño}
\label{sec:proceso_diseno:01}
\pendiente{Texto del apartado pendiente de redactar. Las figuras
\ref{fig:proceso_diseno} y \ref{fig:subproceso_dimensionamiento} ya están
puestas; falta el texto que las explica.}

La figura \ref{fig:proceso_diseno} recoge el proceso que sigue el equipo para
diseñar un componente, desde que se asigna la tarea hasta que el resultado se
entrega al encargado de área. En trazo continuo está el proceso tal y como se
planteó; en trazo discontinuo rojo, los pasos propuestos que todavía no se han
acordado y que se justifican al final del apartado.

Los triángulos numerados son enlaces a subprocesos que se desarrollan por
separado. El único desarrollado hasta ahora es el \textsc{ii}, el
dimensionamiento preliminar, que recoge la figura
\ref{fig:subproceso_dimensionamiento}.

% ---------------------------------------------------------------------
\begin{figure}[p]
\centering
\resizebox{!}{0.90\textheight}{%
\begin{tikzpicture}[
    proceso/.append style={font=\sffamily\footnotesize},
    decision/.append style={font=\sffamily\footnotesize},
    terminal/.append style={font=\sffamily\footnotesize\bfseries,text width=44mm},
    nota/.append style={font=\sffamily\scriptsize},
    conector/.append style={font=\sffamily\scriptsize\bfseries,minimum size=9mm},
    etiq/.append style={font=\sffamily\scriptsize},
    cabecera/.style={proceso,text width=24mm}]

% ---- Flujo principal ----
\node[terminal] (t1)  at (0,0)      {Asignación de tareas};

\node[cabecera] (b1) at (-4.65,-1.75) {Documentos aplicables};
\node[cabecera] (b2) at (-1.55,-1.75) {Objetivos};
\node[cabecera] (b3) at ( 1.55,-1.75) {Método de cálculo};
\node[cabecera,propuesto] (b4) at ( 4.65,-1.75) {Interfaces y restricciones};

\node[proceso]  (p1) at (0,-3.55)  {Dimensionamiento preliminar};
\node[conector] (c2) at (2.7,-3.55) {II};
\node[nota]     (n1) at (6.9,-3.55)
      {\textbf{Pasos}\\[1pt]$\cdot$ Selección de material\\
       $\cdot$ Método de fabricación\\$\cdot$ Ubicación de componentes};

\node[decision] (d1) at (0,-5.35)  {¿Resultado válido?};
\node[proceso,propuesto] (p2) at (0,-7.05) {Registro de decisión};

\node[proceso]  (p3) at (0,-8.75)  {Diseño CAD};
\node[conector] (c3) at (2.7,-8.75) {III};
\node[nota]     (n2) at (6.9,-8.75)
      {\textbf{Pasos}\\[1pt]$\cdot$ Ensamblaje\\$\cdot$ Fabricación\\
       $\cdot$ Material\\$\cdot$ Normativa};

\node[decision] (d2) at (0,-10.55) {¿CAD válido?};

\node[proceso]  (p4) at (0,-12.25) {Simulación MEF};
\node[conector] (c4) at (2.7,-12.25) {IV};
\node[nota]     (n3) at (6.9,-12.25)
      {\textbf{Pasos}\\[1pt]$\cdot$ Condiciones de contorno\\
       $\cdot$ Tipo de vinculación\\$\cdot$ Mallado\\$\cdot$ Convergencia};

\node[decision] (d3) at (0,-14.05) {¿Simulación válida?};
\node[proceso]  (p5) at (0,-15.85) {Revisión de objetivos};
\node[decision] (d4) at (0,-17.65) {¿Objetivos cumplidos?};
\node[nota]     (n4) at (-5.2,-17.65)
      {Revisión completa\\Detección de fallos\\Iteración};

\node[proceso,text width=42mm] (p6) at (0,-19.75)
      {Elaboración de BOM $\cdot$ Informe de diseño $\cdot$ Informe de coste
       $\cdot$ Informe de simulaciones $\cdot$ Hoja de procesos};
\node[conector] (c5) at (4.4,-19.75) {V};

\node[decision,propuesto] (d5) at (0,-21.85) {¿Revisado por otra persona?};
\node[terminal] (t2) at (0,-23.65) {Enviar al encargado de área};
\node[proceso,propuesto,text width=34mm] (p7) at (0,-25.35)
      {Fabricación y validación experimental};

% ---- Flechas del tronco ----
\draw[flujo] (t1.south) -- ++(0,-0.35) coordinate (r1);
\draw[draw=fusalPrim!85,semithick] (b1.north|-r1) -- (b4.north|-r1);
\foreach \n in {b1,b2,b3} \draw[flujo] (\n.north|-r1) -- (\n.north);
\draw[flujoprop] (b4.north|-r1) -- (b4.north);

\coordinate (r2) at (0,-2.65);
\foreach \n in {b1,b2,b3} \draw[draw=fusalPrim!85,semithick] (\n.south) -- (\n.south|-r2);
\draw[draw=fusalSec,dashed,dash pattern=on 2pt off 1.5pt] (b4.south) -- (b4.south|-r2);
\draw[draw=fusalPrim!85,semithick] (b1.south|-r2) -- (b4.south|-r2);
\draw[flujo] (r2) -- (p1.north);

\draw[flujo] (p1) -- (d1);
\draw[flujoprop] (d1) -- node[etiq,pos=0.45]{Sí} (p2);
\draw[flujoprop] (p2) -- (p3);
\draw[flujo] (d2) -- node[etiq,pos=0.45]{Sí} (p4);
\draw[flujo] (p3) -- (d2);
\draw[flujo] (d3) -- node[etiq,pos=0.45]{Sí} (p5);
\draw[flujo] (p4) -- (d3);
\draw[flujo] (p5) -- (d4);
\draw[flujo] (d4) -- node[etiq,pos=0.40]{Sí} (p6);
\draw[flujoprop] (p6) -- (d5);
\draw[flujoprop] (d5) -- node[etiq,pos=0.45]{Sí} (t2);
\draw[flujoprop] (t2) -- (p7);

% ---- Conectores y notas ----
\foreach \p/\c in {p1/c2,p3/c3,p4/c4} \draw[aviso] (\p.east) -- (\c.west);
\draw[aviso] (p6.east) -- (c5.west);
\foreach \c/\n in {c2/n1,c3/n2,c4/n3} \draw[aviso] (\c.east) -- (\n.west);
\draw[aviso] (d4.west) -- node[etiq,pos=0.45]{No} (n4.east);

% ---- Realimentaciones ----
\draw[flujo] (d1.west) -- node[etiq,pos=0.3]{No} ++(-2.1,0) |- (p1.west);
\draw[flujo] (d2.west) -- node[etiq,pos=0.3]{No} ++(-2.6,0) |- ([yshift=-1.2mm]p1.west);
\draw[flujo] (d3.west) -- node[etiq,pos=0.3]{No} ++(-3.1,0) |- ([yshift=1.2mm]p1.west);
\draw[flujoprop] (n4.west) -- ++(-1.1,0) |- ([yshift=-2.4mm]p1.west);
\draw[flujoprop] (d5.east) -- node[etiq,pos=0.35]{No} ++(2.0,0) |- (p6.east);
\draw[flujoprop] (p7.east) -- ++(3.9,0) |- (p5.east);

\end{tikzpicture}}
\caption[Proceso de diseño de un componente]{Proceso de diseño de un
componente. Trazo continuo: el proceso tal y como se planteó. Trazo discontinuo
rojo: pasos propuestos, pendientes de acordar. Los triángulos enlazan con
subprocesos desarrollados aparte.}
\label{fig:proceso_diseno}
\end{figure}

% ---------------------------------------------------------------------
\begin{figure}[p]
\centering
\begin{tikzpicture}[
    proceso/.append style={font=\sffamily\footnotesize},
    decision/.append style={font=\sffamily\footnotesize},
    terminal/.append style={font=\sffamily\footnotesize\bfseries,text width=40mm},
    nota/.append style={font=\sffamily\scriptsize},
    conector/.append style={font=\sffamily\scriptsize\bfseries,minimum size=9mm},
    etiq/.append style={font=\sffamily\scriptsize}]

\node[conector] (ci) at (0,0.9) {II};
\node[proceso]  (s1) at (0,-0.5)  {Búsqueda de referencias};
\node[nota]     (m1) at (4.6,-0.5) {$\cdot$ Libros\\$\cdot$ TFG y TFM\\
       $\cdot$ Foros\\$\cdot$ Vídeos y documentación de otros equipos};

\node[decision] (e1) at (0,-2.4)  {¿Método fiable?};
\node[proceso]  (s2) at (0,-4.3)  {Desarrollo analítico};
\node[proceso]  (s3) at (0,-6.2)  {Programar el cálculo};
\node[nota]     (m2) at (-4.6,-6.2) {$\cdot$ Excel\\$\cdot$ MathCAD\\
       $\cdot$ MATLAB\\$\cdot$ Programación};

\node[proceso]  (s4) at (0,-8.1)  {Validación y comparación};
\node[nota]     (m3) at (4.6,-8.1) {$\cdot$ Objetivos\\
       $\cdot$ Referencias anteriores\\$\cdot$ Ensayo, cuando exista};

\node[decision] (e2) at (0,-10.0) {¿Válido?};
\node[proceso,propuesto,text width=34mm] (s5) at (0,-11.9)
      {Documentar hipótesis y archivar el modelo};
\node[terminal] (fin) at (0,-13.7) {Continuar: diseño CAD};

\draw[flujo] (ci) -- (s1);
\draw[flujo] (s1) -- (e1);
\draw[flujo] (e1) -- node[etiq,pos=0.45]{Sí} (s2);
\draw[flujo] (s2) -- (s3);
\draw[flujo] (s3) -- (s4);
\draw[flujo] (s4) -- (e2);
\draw[flujoprop] (e2) -- node[etiq,pos=0.45]{Sí} (s5);
\draw[flujoprop] (s5) -- (fin);

\draw[flujo] (e1.east) -- node[etiq,pos=0.35]{No} ++(2.2,0) |- (s3.east);
\draw[flujo] (e2.west) -- node[etiq,pos=0.25]{No} ++(-5.4,0) |- (s1.west);

\draw[aviso] (s1.east) -- (m1.west);
\draw[aviso] (s3.west) -- (m2.east);
\draw[aviso] (s4.east) -- (m3.west);

\end{tikzpicture}
\caption[Subproceso \textsc{ii}: dimensionamiento preliminar]{Subproceso
\textsc{ii}: dimensionamiento preliminar. Los subprocesos \textsc{iii} (diseño
CAD), \textsc{iv} (simulación MEF) y \textsc{v} (documentación) están pendientes
de desarrollar.}
\label{fig:subproceso_dimensionamiento}
\end{figure}

\subsection{Qué falta en el proceso}

Al revisar el diagrama aparecen seis huecos. Cuatro de ellos ---los números
\ref{item:huecos:revision}, \ref{item:huecos:ddr}, \ref{item:huecos:validacion} y
\ref{item:huecos:interfaces}--- están dibujados en rojo en la figura
\ref{fig:proceso_diseno}. Los otros dos no, porque antes hay que decidir dónde
encajan.

\begin{enumerate}
  \item \textbf{No hay salida de los bucles.} Las tres preguntas de validez
        devuelven al dimensionamiento preliminar sin límite. Hace falta un
        criterio de escalado: tras dos iteraciones sin converger, el problema
        pasa al encargado de área. Si no, una tarea puede consumir media
        temporada sin que nadie se entere.
  \item\label{item:huecos:revision} \textbf{Nadie revisa el trabajo antes de
        entregarlo.} El capítulo
        \ref{ch:herramientas_cad_datos} exige que una pieza no se libere sin que
        la revise alguien distinto del autor. Esa comprobación tiene que estar en
        el proceso, no en la buena voluntad de cada uno.
  \item\label{item:huecos:ddr} \textbf{El registro de decisión no aparece.} El diagrama recoge qué se
        hace, pero no deja rastro de \emph{por qué} se eligió una alternativa y
        no otra. Sin eso, el año que viene se vuelve a empezar --- que es
        justamente lo que este libro intenta evitar.
  \item \textbf{El coste entra demasiado tarde.} Aparece solo en la
        documentación final, cuando ya no se puede cambiar nada. Debería ser un
        criterio de la primera pregunta de validez, junto con la fabricabilidad.
  \item\label{item:huecos:validacion} \textbf{No hay realimentación desde la fabricación ni desde la pista.}
        El proceso termina al entregar los documentos. Lo que pasa después ---
        que la pieza no se pueda fabricar como estaba prevista, o que rompa ---
        no vuelve a entrar en ningún sitio.
  \item\label{item:huecos:interfaces} \textbf{Las interfaces con otros subsistemas no se acuerdan.} El bloque
        de partida recoge documentos, objetivos y método, pero no con quién hay
        que hablar antes de fijar una cota compartida (capítulo
        \ref{ch:coche_como_sistema}).
\end{enumerate}

\pendiente{Decidir con el equipo cuáles de estos seis pasos se incorporan, y
actualizar la figura. Lo que se acuerde deja de ser trazo discontinuo.}

\pendiente{Transcripción de la versión manuscrita: revisar las etiquetas
«Desarrollo analítico» y «Validación y comparación», y el contenido del recuadro
de herramientas de cálculo, que se han leído con dudas.}

\pendiente{El boceto original incluye un recuadro hexagonal, enlazado al conector
\textsc{i}, que dice «Aclaración: Coordinador, Encargado y Supervisor». No se ha
llevado a la figura porque no está claro a qué paso del proceso se refiere.
Decidir si es la definición de los tres papeles que intervienen ---y entonces va
en el apartado \ref{sec:proceso_diseno:04}, junto al registro de decisiones--- o
si marca los puntos del flujo en los que cada uno interviene, en cuyo caso hay
que dibujarlo.}

\section{Generación de conceptos}
\label{sec:proceso_diseno:02}
\pendiente{Pendiente de redactar.}

\section{Matrices de decisión}
\label{sec:proceso_diseno:03}
\pendiente{Pendiente de redactar.}

\section{Registro de decisiones}
\label{sec:proceso_diseno:04}
\pendiente{Pendiente de redactar.}

\section{Congelación del diseño}
\label{sec:proceso_diseno:05}
\pendiente{Pendiente de redactar.}

% ---------------------------------------------------------------------
%  Cierre del capítulo: los cuatro recuadros son obligatorios.
% ---------------------------------------------------------------------

\begin{caso}
\redactar{Qué hicimos nosotros: decisión tomada, datos de partida, resultado en pista y qué haríamos distinto.}
\end{caso}

\begin{ojo}
\begin{itemize}
  \item \redactar{Error que repetimos cada temporada en este tema.}
\end{itemize}
\end{ojo}

\begin{entregable}
\redactar{Qué produce el lector al terminar: plantilla, hoja de cálculo, checklist o apartado del informe.}
\end{entregable}

\begin{juez}
\begin{enumerate}
  \item \redactar{Pregunta tipo del design event sobre este capítulo.}
\end{enumerate}
\end{juez}

% !TeX root = ../../main.tex
% =====================================================================
%  Herramientas de trabajo: CAD, datos y versiones
%  Responsable:
%  Última revisión:
% =====================================================================
\chapter{Herramientas de trabajo: CAD, datos y versiones}
\label{ch:herramientas_cad_datos}

\begin{objetivos}
  \item Montar el ensamblaje maestro del coche y explicar por qué ningún
        subsistema referencia directamente a otro.
  \item Definir el sistema de coordenadas del coche y la transformación al
        convenio de dinámica vehicular.
  \item Asignar correctamente un código a cualquier pieza, conjunto, componente
        comercial o utillaje del coche.
  \item Distinguir cuándo un cambio es una revisión y cuándo es una pieza nueva.
  \item Preparar un paquete de fabricación que un taller externo pueda ejecutar
        sin llamarte por teléfono.
\end{objetivos}

\fichasesion{90 minutos}{Capítulo \ref{ch:coche_como_sistema}}%
{Ordenador con el CAD del equipo instalado y acceso al servidor}

Este capítulo no va de aprender a usar el CAD. Va de las convenciones que hacen
que el trabajo de quince personas durante nueve meses siga siendo utilizable
cuando ninguna de esas quince siga en el equipo. Es, junto con el capítulo
\ref{ch:documentacion_traspaso}, el que más directamente sirve al propósito de
este libro.

\subsection*{Nuestras herramientas}

El equipo diseña en \textbf{Fusion 360}, con la intención de pasar a
\textbf{SolidWorks} más adelante. Eso condiciona tres cosas de este capítulo:

\begin{itemize}
  \item \textbf{Las convenciones de este capítulo son independientes de la
        herramienta, y eso es deliberado.} Los códigos viven en el campo
        \emph{Part Number} del componente y en el nombre del fichero, no en una
        función propia de un programa. Si un día se cambia de CAD, los códigos, la
        estructura de carpetas y el registro de piezas siguen valiendo.
  \item \textbf{Fusion resuelve por defecto parte del apartado
        \ref{sec:herramientas_cad_datos:04}}: guarda el histórico de versiones de
        cada fichero y centraliza los datos en la nube. Lo que \emph{no} resuelve
        solo es la organización: los permisos por carpeta, quién es responsable de
        qué y las congelaciones de hitos hay que montarlos igual.
  \item \textbf{La migración no es gratis.} Al pasar geometría de un CAD a otro
        por un formato neutro tipo STEP viaja el sólido, pero no el árbol de
        operaciones: en el programa de destino aparece una pieza que no se puede
        editar paramétricamente. Por eso una migración se hace al empezar un coche
        nuevo, remodelando, y no a mitad de temporada.
\end{itemize}

\pendiente{Cuando se cierre la decisión sobre SolidWorks, añadir aquí las
instrucciones concretas de cada programa: plantillas de pieza y de plano, cómo se
rellenan las propiedades y cómo se exporta la lista de materiales.}

% ---------------------------------------------------------------------
\section{Estructura del ensamblaje maestro}
\label{sec:herramientas_cad_datos:01}
\index{ensamblaje maestro}

Todo el coche cuelga de un único ensamblaje de nivel superior, y por debajo hay
exactamente un nivel de conjuntos de subsistema:

\begin{center}
\begin{minipage}{0.8\textwidth}
\ttfamily\small
FS26-GEN-A000\quad coche completo\\
\hspace*{1em}FS26-GEN-S001\quad esqueleto (referencia de todos)\\
\hspace*{1em}FS26-CHS-A000\quad conjunto chasis\\
\hspace*{1em}FS26-SUS-A000\quad conjunto suspensión\\
\hspace*{2em}FS26-SUS-A001\quad esquina delantera izquierda\\
\hspace*{2em}FS26-SUS-A002\quad esquina delantera derecha\\
\hspace*{1em}FS26-AER-A000\quad conjunto aerodinámica\\
\hspace*{1em}\ldots
\end{minipage}
\end{center}

Y una única regla, que es la más importante de todo el capítulo:

\begin{clave}
\textbf{Ningún subsistema referencia geometría de otro subsistema.} Todas las
referencias externas van contra el esqueleto, y solo contra el esqueleto. Si la
mangueta necesita saber dónde está el anclaje del trapecio, ese punto está en el
esqueleto; no se coge del chasis.
\end{clave}

El motivo es práctico: en cuanto un modelo referencia a otro, abrir una pieza
obliga a cargar medio coche, un cambio en el chasis puede romper silenciosamente
la suspensión, y dos personas no pueden trabajar a la vez. Con el esqueleto como
único punto de contacto, cada subsistema es independiente y los cambios que
afectan a todos pasan por un sitio en el que se ven.

% ---------------------------------------------------------------------
\section{Esqueleto y puntos duros}
\label{sec:herramientas_cad_datos:02}
\index{esqueleto}\index{puntos duros}

El esqueleto (\texttt{FS26-GEN-S001}) es una pieza sin volumen que contiene solo
geometría de referencia. Es el contrato entre subsistemas hecho fichero.

\subsection{Sistema de coordenadas}

Antes que nada hay que fijar el origen y la orientación de los ejes, escribirlo,
y no volver a discutirlo:

\begin{itemize}
  \item \textbf{Origen}: intersección del plano de simetría del coche, el plano
        del suelo y el eje delantero (centro de las ruedas delanteras proyectado
        al suelo), con el coche a la altura de marcha de diseño.
  \item \textbf{$X$} positivo hacia atrás. Así la batalla y las posiciones
        longitudinales son números positivos crecientes hacia la cola.
  \item \textbf{$Y$} positivo hacia la derecha del piloto.
  \item \textbf{$Z$} positivo hacia arriba.
\end{itemize}

\begin{ojo}
El convenio de ejes del CAD \textbf{no} coincide con el convenio SAE que se usa
en dinámica vehicular (véase la nomenclatura), donde $Z$ apunta hacia abajo, lo
que provoca errores de signo al pasar cargas de un sitio a otro.
La transformación entre ambos se escribe \emph{una vez}, en la nomenclatura de
este libro, y se aplica siempre desde ahí. Nunca «a ojo».
\end{ojo}

\subsection{Contenido del esqueleto}

\begin{center}
\begin{tabularx}{\textwidth}{@{}lX@{}}
\toprule
\textbf{Elemento} & \textbf{Qué incluye} \\
\midrule
Planos de referencia & Simetría, suelo, eje delantero, eje trasero \\
Cotas maestras & Batalla, vías delantera y trasera, altura de marcha de diseño \\
Centros de rueda & Los cuatro, en posición de diseño \\
Puntos duros de suspensión & Todos los anclajes, chasis y mangueta (capítulo \ref{ch:puntos_duros}) \\
Referencia del piloto & Punto del asiento, plantillas de reglamento \\
Volúmenes reservados & Espacios que un subsistema no puede invadir \\
\bottomrule
\end{tabularx}
\end{center}

Los volúmenes reservados merecen una mención: son cajas sin más significado que
«aquí no entra nadie más». Reservar desde el principio el volumen del radiador,
el del acumulador o el del recorrido del mazo evita la mitad de los conflictos
de integración del capítulo \ref{ch:coche_como_sistema}.

\subsection{Quién puede tocarlo}

El esqueleto tiene \textbf{un único responsable} con permiso de escritura ---
normalmente el responsable de integración o el de chasis. Cualquier cambio en él
afecta a todo el coche, así que se comunica al equipo entero y se anota con
fecha y motivo. Un cambio de punto duro que no se comunica se acaba descubriendo
en el montaje, con las piezas ya fabricadas.

% ---------------------------------------------------------------------
\section{Nomenclatura y numeración de piezas}
\label{sec:herramientas_cad_datos:03}
\index{nomenclatura}\index{codificación de piezas}

\subsection{Anatomía del código}

Toda pieza, conjunto, componente comprado y utillaje del coche tiene un código
con esta forma:

\begin{figure}[H]
\centering
\begin{tikzpicture}[
    font=\sffamily\footnotesize,
    bloque/.style={anchor=west,inner xsep=0pt,inner ysep=2.5pt,outer sep=0pt,
                   font=\ttfamily\LARGE},
    etiq/.style={anchor=north,align=center,text width=2.2cm,fusalPrim},
    guia/.style={fusalGris,thin}]
  % El código, encadenado para que no haya huecos entre grupos
  \node[bloque] (g1) at (0,0)      {FS26};
  \node[bloque] (g2) at (g1.east)  {-};
  \node[bloque] (g3) at (g2.east)  {SUS};
  \node[bloque] (g4) at (g3.east)  {-};
  \node[bloque] (g5) at (g4.east)  {P};
  \node[bloque] (g6) at (g5.east)  {014};

  % Etiquetas, repartidas para que no se solapen
  \node[etiq] (e1) at (-0.9,-1.5) {\textbf{Proyecto}\\coche y temporada};
  \node[etiq] (e2) at ( 1.4,-1.5) {\textbf{Subsistema}\\3 letras, lista cerrada};
  \node[etiq] (e3) at ( 3.7,-1.5) {\textbf{Tipo}\\1 letra};
  \node[etiq] (e4) at ( 6.0,-1.5) {\textbf{Secuencial}\\3 dígitos};

  \draw[guia] (g1.south) -- ++(0,-0.35) -- (e1.north);
  \draw[guia] (g3.south) -- ++(0,-0.35) -- (e2.north);
  \draw[guia] (g5.south) -- ++(0,-0.35) -- (e3.north);
  \draw[guia] (g6.south) -- ++(0,-0.35) -- (e4.north);
\end{tikzpicture}
\caption{Anatomía del código de pieza. Trece caracteres, legibles por teléfono y
ordenables en el explorador de archivos.}
\label{fig:anatomia_codigo}
\end{figure}

\begin{clave}
El código \textbf{identifica}, no describe. Nunca se codifica en él nada que
pueda cambiar sin que cambie la pieza: material, proveedor, masa, autor o fecha.
Todo eso vive en las propiedades del modelo, donde se puede corregir sin
renumerar medio coche.
\end{clave}

\subsection{Códigos de subsistema}

Lista cerrada. Si hace falta uno nuevo, se añade aquí y se comunica; no se
improvisa.

\begin{center}
\begin{tabularx}{\textwidth}{@{}llX@{}}
\toprule
\textbf{Código} & \textbf{Subsistema} & \textbf{Alcance} \\
\midrule
\ttfamily GEN & General        & Esqueleto, conjunto general, tornillería propia, misceláneo \\
\ttfamily CHS & Chasis         & Estructura primaria, anclajes, atenuador, mamparos \\
\ttfamily SUS & Suspensión     & Trapecios, manguetas, bujes, balancines, tirantes, muelles, amortiguadores, barras \\
\ttfamily DIR & Dirección      & Cremallera, columna, volante, bieletas \\
\ttfamily FRE & Frenos         & Discos, pinzas, bombas, pedalera de freno, conductos \\
\ttfamily RUE & Ruedas         & Llantas, neumáticos, tuercas y centros de rueda \\
\ttfamily AER & Aerodinámica   & Alerones, fondo, difusor, carrocería, soportes \\
\ttfamily REF & Refrigeración  & Radiadores, ventiladores, conductos, bomba, circuito \\
\ttfamily TRA & Transmisión    & Reductora, diferencial, palieres, cadena o correa, resguardos \\
\ttfamily MOT & Motor          & Grupo motopropulsor y sus soportes \\
\ttfamily ACU & Acumulador     & Celdas, contenedor, refrigeración de batería (vehículos eléctricos) \\
\ttfamily ELE & Electrónica    & Placas, cajas, sensores, mazos \\
\ttfamily ERG & Ergonomía      & Asiento, pedalera, reposacabezas, arneses, cockpit \\
\bottomrule
\end{tabularx}
\end{center}

\pendiente{Revisar esta lista con los responsables de subsistema y cerrarla antes
de empezar a modelar. Cambiarla a mitad de temporada cuesta mucho más de lo que parece.}

\subsection{Letras de tipo}

\begin{center}
\begin{tabularx}{\textwidth}{@{}clX@{}}
\toprule
\textbf{Letra} & \textbf{Tipo} & \textbf{Uso} \\
\midrule
\ttfamily A & Conjunto & Ensamblajes. El \ttfamily A000 \normalfont de cada subsistema es su conjunto principal \\
\ttfamily P & Pieza    & Piezas que fabricamos nosotros o encargamos fabricar según plano \\
\ttfamily C & Comercial & Componentes comprados de catálogo \\
\ttfamily U & Utillaje & Moldes, plantillas, útiles de soldadura, bancos de ensayo \\
\ttfamily S & Simulación & Geometría de referencia, maquetas y volúmenes reservados. No se fabrica \\
\bottomrule
\end{tabularx}
\end{center}

El secuencial son tres dígitos, empieza en \texttt{001} y es independiente para
cada combinación de subsistema y tipo. Es decir, pueden coexistir
\texttt{FS26-SUS-P014} y \texttt{FS26-SUS-C014}: son cosas distintas y no se
confunden porque la letra los separa.

\subsection{Las seis reglas de numeración}

\begin{enumerate}
  \item \textbf{Un número no se reutiliza nunca}, ni aunque la pieza se abandone.
        Reutilizarlo hace que dos documentos distintos digan cosas distintas del
        mismo código.
  \item \textbf{Izquierda y derecha son dos códigos distintos.} Aunque sean
        espejo, son dos piezas que se fabrican, se compran y se cuestan por
        separado, y así deben aparecer en el BOM. La relación de simetría se anota
        en las propiedades.
  \item \textbf{Los componentes comerciales también llevan código nuestro}, del
        tipo \texttt{C}. La referencia del fabricante va en las propiedades: así,
        cambiar de proveedor no obliga a tocar el diseño.
  \item \textbf{Si la pieza nueva es intercambiable con la vieja, es una
        revisión.} Si no lo es, es un código nuevo. Esta es la única prueba que
        hay que aplicar, y se aplica preguntando: ¿puedo montar la nueva donde
        estaba la vieja y que el coche funcione igual?
  \item \textbf{El código se asigna al crear el fichero}, no al terminarlo. Una
        pieza sin código no existe y no se puede referenciar.
  \item \textbf{El registro de códigos es único y compartido.} Una hoja con una
        fila por código, y quien asigna el siguiente lo anota en el momento. Sin
        ese registro aparecen números duplicados en cuanto trabajan dos personas
        a la vez.
\end{enumerate}

\subsection{Nombres de fichero}

El nombre del fichero es el código, un guion bajo, y una descripción corta:

\begin{center}
\ttfamily\large FS26-SUS-P014\_trapecio-del-sup-izq.\textrm{\itshape ext}
\end{center}

Reglas de la parte descriptiva:

\begin{itemize}
  \item Minúsculas, sin acentos, sin eñes, sin espacios. Separador: guion.
  \item Corta: unos 30 caracteres. Es una etiqueta, no una frase.
  \item Abreviaturas fijas para la posición: \texttt{del}/\texttt{tra},
        \texttt{sup}/\texttt{inf}, \texttt{izq}/\texttt{der}, en ese orden.
\end{itemize}

\begin{ojo}
Lo que \textbf{nunca} aparece en el nombre de un fichero:

\begin{center}
\ttfamily\small
trapecio\_v2.ext \quad trapecio\_FINAL.ext \quad trapecio\_final\_bueno.ext\\
trapecio\_2026-03-14.ext \quad trapecio\_JL.ext \quad trapecio\_OK\_este\_si.ext
\end{center}

\noindent
La versión, la fecha y el autor los guarda el sistema de archivos y el control de
versiones, y los guarda bien. Escritos en el nombre, lo único que consiguen es que
dentro de unos meses nadie sepa cuál es el fichero bueno.
\end{ojo}

\subsection{Revisiones y estados}
\index{revisiones}

La revisión es una letra: \texttt{A} para la primera versión liberada,
\texttt{B} para la siguiente, y así. \textbf{Solo hay revisiones después de la
primera liberación}: mientras la pieza se está diseñando no se revisa, se
modifica.

Cada pieza tiene además un estado, y el estado decide qué se puede hacer con ella:

\begin{center}
\begin{tabularx}{\textwidth}{@{}lXl@{}}
\toprule
\textbf{Estado} & \textbf{Significado} & \textbf{¿Se fabrica?} \\
\midrule
En diseño   & Trabajo en curso. Puede cambiar en cualquier momento & No \\
En revisión & Terminada, pendiente de que la revise otra persona & No \\
Liberada    & Revisada y aprobada. Plano firmado & \textbf{Sí} \\
Obsoleta    & Sustituida o abandonada. Se conserva, no se borra & No \\
\bottomrule
\end{tabularx}
\end{center}

\begin{clave}
Al taller solo va material \textbf{liberado}, y liberar exige que lo haya revisado
alguien que no sea el autor. Cuesta unos minutos y detecta errores que el autor ya
no ve.
\end{clave}

Cuando una pieza liberada cambia, la revisión sube una letra, se anota en la
tabla de revisiones del plano (qué cambió y por qué) y \emph{se avisa a quien
tenga la revisión anterior en la mano}: si no, seguirá trabajando con un plano
obsoleto sin saberlo.

\subsection{Propiedades obligatorias del modelo}
\index{propiedades del modelo}

Cada pieza lleva rellenas estas propiedades personalizadas. No es burocracia: son
las que alimentan automáticamente el cajetín del plano, la lista de materiales y
el informe de costes del capítulo \ref{ch:evento_costes}. Rellenarlas al crear la
pieza cuesta un minuto; reconstruirlas en junio para el informe de costes cuesta
semanas.

\begin{center}
\begin{tabularx}{\textwidth}{@{}lXl@{}}
\toprule
\textbf{Propiedad} & \textbf{Contenido} & \textbf{Obligatoria} \\
\midrule
\ttfamily CODIGO        & Código completo de la pieza. En Fusion, el campo
                          \emph{Part Number} del componente & Sí \\
\ttfamily DENOMINACION  & Nombre en castellano, como aparece en el plano.
                          En Fusion, \emph{Description} & Sí \\
\ttfamily ESTADO        & En diseño / en revisión / liberada / obsoleta & Sí \\
\ttfamily REVISION      & Letra de revisión & Sí \\
\ttfamily AUTOR         & Quien la diseña & Sí \\
\ttfamily REVISOR       & Quien la aprueba al liberarla & Al liberar \\
\ttfamily FECHA         & Fecha de liberación & Al liberar \\
\ttfamily MATERIAL      & Designación normalizada, no «aluminio» & Sí \\
\ttfamily MASA          & Calculada por el CAD & Automática \\
\ttfamily PROCESO       & Mecanizado, soldadura, laminado, impresión\ldots & Sí \\
\ttfamily CANTIDAD      & Unidades por coche & Sí \\
\ttfamily PROVEEDOR     & Solo para comerciales y subcontratados & Si aplica \\
\ttfamily REF\_PROVEEDOR & Referencia de catálogo del fabricante & Si aplica \\
\ttfamily TRATAMIENTO   & Térmico o superficial, si lo lleva & Si aplica \\
\ttfamily SUBSISTEMA    & Para agrupar el BOM & Sí \\
\bottomrule
\end{tabularx}
\end{center}

\subsection{Carpetas}

\begin{center}
\begin{minipage}{0.8\textwidth}
\ttfamily\small
FS26\_CAD/\\
\hspace*{1em}00\_GEN/\quad esqueleto, conjunto general, biblioteca de tornillería\\
\hspace*{1em}01\_CHS/\\
\hspace*{2em}Piezas/\\
\hspace*{2em}Conjuntos/\\
\hspace*{2em}Comerciales/\\
\hspace*{2em}Planos/\\
\hspace*{2em}Fabricacion/\quad paquetes listos para taller (PDF + STEP + DXF)\\
\hspace*{1em}02\_SUS/\quad\textrm{(misma estructura)}\\
\hspace*{1em}\ldots\\
\hspace*{1em}98\_Obsoleto/\quad nada se borra, todo se mueve aquí\\
\hspace*{1em}99\_Congelaciones/\quad copias de solo lectura de cada hito
\end{minipage}
\end{center}

% ---------------------------------------------------------------------
\section{Control de versiones y copias de seguridad}
\label{sec:herramientas_cad_datos:04}
\index{copias de seguridad}

El CAD de un coche es el activo más valioso que produce el equipo y el más fácil
de perder. Las reglas son pocas y no admiten excepción:

\begin{enumerate}
  \item \textbf{Hay un único original, y está en el servidor del equipo.} Lo que
        hay en un portátil es una copia de trabajo, nunca la referencia. Un
        proyecto que vive en el ordenador de alguien se pierde el día que esa
        persona se va, y se va siempre.
  \item \textbf{Una carpeta, un responsable con permiso de escritura.} Los demás
        leen. Si alguien necesita modificar una pieza que no es suya, se la pide
        a su responsable.
  \item \textbf{Nadie abre el mismo fichero desde dos sitios.} Si el CAD no tiene
        gestor de datos que bloquee ficheros, se suple con una hoja compartida de
        control: quién tiene qué abierto y desde cuándo. Es rudimentario y funciona.
  \item \textbf{Copia de seguridad automática diaria}, conservando al menos
        treinta días. Automática: lo que depende de que alguien se acuerde no es
        una copia de seguridad.
  \item \textbf{Congelación en cada hito.} Al cerrar el diseño de un subsistema se
        guarda una copia completa de \emph{solo lectura} en
        \texttt{99\_Congelaciones/}, con el nombre
        \texttt{FS26\_congelacion\_2026-02-15}. Es la única forma de poder
        contestar dentro de un año a «¿qué mandamos a fabricar exactamente?».
  \item \textbf{Al terminar la temporada, archivo completo del coche tal y como
        se fabricó}, incluyendo planos, paquetes de fabricación y el registro de
        códigos. Ese archivo es el que se consultará dentro de tres años, y por
        entonces no quedará nadie a quien preguntar.
\end{enumerate}

\begin{ojo}
\begin{itemize}
  \item Trabajar con el CAD en una carpeta sincronizada a la nube y abierta desde
        dos ordenadores a la vez. Las herramientas de sincronización resuelven los
        conflictos duplicando ficheros, y en un ensamblaje eso rompe las
        referencias sin avisar.
  \item Confundir sincronización con copia de seguridad. Si borras una carpeta,
        la nube sincroniza el borrado obedientemente.
  \item Guardar las congelaciones con permiso de escritura. Acaban modificadas.
\end{itemize}
\end{ojo}

% ---------------------------------------------------------------------
\section{Planos, acotación y entrega a taller}
\label{sec:herramientas_cad_datos:05}
\index{planos}

\subsection{El plano es el contrato}

Lo que no está en el plano no existe. El taller no sabe lo que tenías en la
cabeza, no va a abrir tu modelo y, si algo es ambiguo, elegirá la interpretación
más barata de fabricar --- que es lo correcto por su parte.

\begin{itemize}
  \item \textbf{Formato A3} por defecto, con cajetín rellenado automáticamente
        desde las propiedades del modelo.
  \item \textbf{Acotación funcional}: se acota desde la superficie que define la
        función de la pieza, no desde la esquina que resultaba cómoda. Las
        referencias de acotación deben coincidir con las de montaje.
  \item \textbf{Tolerancia general} indicada en el cajetín (por ejemplo
        ISO 2768-mK) y tolerancias estrechas \emph{solo} donde importan. Cada
        tolerancia apretada es dinero y tiempo de máquina; si están todas
        apretadas, en realidad no has decidido ninguna.
  \item \textbf{Acabado superficial} solo donde tenga función.
  \item \textbf{Tabla de revisiones} en el plano, con fecha y motivo de cada una.
\end{itemize}

\subsection{Paquete de fabricación}

Cada pieza que va a taller se entrega como un paquete completo en
\texttt{Fabricacion/}, con el código y la revisión en el nombre:

\begin{center}
\begin{tabularx}{\textwidth}{@{}lX@{}}
\toprule
\textbf{Fichero} & \textbf{Para qué} \\
\midrule
\ttfamily \ldots\_planoB.pdf  & El plano firmado. Es el documento contractual \\
\ttfamily \ldots\_revB.step   & Geometría neutra para programar la máquina \\
\ttfamily \ldots\_revB.dxf    & Solo para corte plano: láser, agua, plegado \\
\bottomrule
\end{tabularx}
\end{center}

\begin{clave}
A un proveedor externo se le envía \textbf{siempre} STEP y PDF, nunca el fichero
nativo del CAD. El nativo lleva el histórico de operaciones, las relaciones con el
resto del coche y a menudo información que no queremos repartir; además obliga al
proveedor a tener nuestro mismo programa y versión.
\end{clave}

% ---------------------------------------------------------------------
%  Cierre del capítulo
% ---------------------------------------------------------------------

\begin{caso}
\redactar{Rellenar con la situación real: qué CAD usamos, dónde vive el modelo
hoy, si existe esqueleto, qué convenio de códigos veníamos usando y qué se ha
perdido de temporadas anteriores. Conviene ser explícito sobre lo que no
tenemos: es lo que justifica adoptar este capítulo.}
\end{caso}

\begin{ojo}
\begin{itemize}
  \item Empezar a modelar antes de tener el esqueleto. Se acaba rehaciendo todo.
  \item Referenciar geometría de otro subsistema «solo esta vez».
  \item Dejar las propiedades del modelo para el final: en junio, cuando hay que
        entregar el informe de costes, ya no hay quien las reconstruya.
  \item Numerar por orden de fabricación en lugar de por subsistema y tipo.
  \item Enviar a taller un plano de una pieza que no está liberada.
  \item Tolerar en el plano cotas que no se pueden medir con lo que hay en el
        equipo. Si no la puedes verificar, no la puedes exigir.
\end{itemize}
\end{ojo}

\begin{entregable}
La \textbf{guía de estándares CAD del equipo}, que es este capítulo con las
decisiones concretas tomadas:

\begin{enumerate}
  \item Tabla de códigos de subsistema cerrada y aprobada por los responsables.
  \item Fichero de esqueleto creado, con el sistema de coordenadas documentado y
        la transformación al convenio SAE escrita.
  \item Plantilla de pieza y de plano con las propiedades obligatorias y el
        cajetín ya enlazado.
  \item Estructura de carpetas creada en el servidor, con permisos asignados.
  \item Registro de códigos abierto, con un responsable de asignación.
\end{enumerate}
\end{entregable}

\begin{juez}
\begin{enumerate}
  \item ¿Cómo os aseguráis de que el modelo que está en el taller es el mismo que
        el que habéis analizado?
  \item Enséñame el plano de esta pieza. ¿Por qué está acotada desde ahí?
  \item ¿Por qué esta tolerancia es tan estrecha? ¿Cómo la verificáis?
  \item ¿Qué pasa cuando cambiáis un punto duro de la suspensión? ¿Quién se entera?
  \item Si tu compañero de subsistema deja el equipo mañana, ¿podrías seguir su
        trabajo?
\end{enumerate}
\end{juez}

% !TeX root = ../../main.tex
% =====================================================================
%  Fundamentos de mecánica aplicada
%  Responsable:
%  Última revisión:
% =====================================================================
\chapter{Fundamentos de mecánica aplicada}
\label{ch:mecanica_aplicada}

\begin{objetivos}
  \item Plantear el diagrama de sólido libre de una esquina del coche y sacar de
        él las cargas de cada barra.
  \item Identificar si una pieza está limitada por resistencia, por rigidez o por
        pandeo, y dimensionarla en consecuencia.
  \item Saber cuándo hay que comprobar fatiga y cuándo no merece la pena.
  \item Calcular el par de apriete de un tornillo a partir de la precarga que se
        quiere conseguir, y entender por qué la precarga es lo único que importa.
  \item Elegir material con criterio, sabiendo qué diferencias entre materiales
        son reales y cuáles son mitología.
\end{objetivos}

\fichasesion{90 minutos}{Resistencia de materiales de primer ciclo}%
{Calculadora, hoja de cálculo, un tornillo y una llave dinamométrica}

Este capítulo no sustituye a una asignatura de resistencia de materiales: da por
sabido lo básico y se centra en las cinco cosas que un equipo de Formula Student
usa constantemente y suele hacer mal. Para profundizar,
\textcite{budynas_shigley_2020} cubre todo lo que hay aquí con mucho más detalle.

% ---------------------------------------------------------------------
\section{Diagramas de sólido libre}
\label{sec:mecanica_aplicada:01}
\index{diagrama de sólido libre}

Todo cálculo estructural empieza aislando un trozo del coche y dibujando todo lo
que actúa sobre él. Parece elemental, pero conviene hacerlo con cuidado: un
diagrama mal planteado da un resultado coherente consigo mismo, y por eso el
error no se detecta después.

El método no tiene misterio:

\begin{enumerate}
  \item \textbf{Aísla} el cuerpo y dibújalo separado de todo lo demás.
  \item \textbf{Sustituye} cada cosa que has quitado por las fuerzas y momentos
        que ejercía.
  \item \textbf{Aplica} equilibrio: $\sum \vect{F} = \vect{0}$ y
        $\sum \vect{M} = \vect{0}$ respecto de un punto cualquiera.
  \item \textbf{Comprueba}: la suma de reacciones debe dar la carga aplicada.
\end{enumerate}

\subsection{El caso canónico: una esquina del coche}

El ejercicio que hay que saber hacer de memoria es este: dada una fuerza en la
huella del neumático, obtener la carga axial de cada barra de la suspensión.

La simplificación que lo hace tratable es que \textbf{los trapecios, el tirante
de dirección y el push\-rod son barras biarticuladas}: rotuladas en ambos
extremos, no pueden transmitir momento, así que solo trabajan a tracción o
compresión pura. Por eso son tubos y por eso el cálculo se reduce a resolver un
sistema de equilibrio de fuerzas concurrentes.

La fuerza de entrada actúa en la \emph{huella de contacto}, no en el centro de
rueda. La diferencia es el radio del neumático multiplicado por la fuerza
lateral, y es un momento que hay que llevar hasta la mangueta.

\begin{ojo}
\begin{itemize}
  \item Aplicar la carga en el centro de rueda en vez de en la huella.
  \item Olvidar el equilibrio de momentos: con solo fuerzas, el sistema parece
        resuelto y no lo está.
  \item Suponer que una barra con casquillos rígidos (no rotulados) trabaja solo
        a axial. Si el extremo no gira libremente, hay flexión.
  \item No comprobar el resultado: si las reacciones no suman la carga aplicada,
        hay un error. Esta comprobación cuesta treinta segundos.
\end{itemize}
\end{ojo}

% ---------------------------------------------------------------------
\section{Tensión, rigidez y pandeo}
\label{sec:mecanica_aplicada:02}
\index{pandeo}

\subsection{Las tres preguntas}

Antes de dimensionar cualquier pieza hay que saber qué la limita, porque cada
límite se combate de forma distinta y a veces con soluciones opuestas:

\begin{center}
\begin{tabularx}{\textwidth}{@{}lXl@{}}
\toprule
\textbf{Limitada por} & \textbf{Pregunta} & \textbf{Qué la mejora} \\
\midrule
Resistencia & ¿Rompe o plastifica? & Material más resistente, más sección \\
Rigidez     & ¿Se deforma demasiado? & Más momento de inercia, mejor forma \\
Pandeo      & ¿Colapsa a compresión? & Más diámetro, menos longitud libre \\
Fatiga      & ¿Rompe tras muchos ciclos? & Radios, acabado, menos concentración \\
\bottomrule
\end{tabularx}
\end{center}

En un coche de Formula Student, \textbf{la mayoría de las piezas estructurales no
están limitadas por resistencia}. Los trapecios los limita el pandeo, el chasis
lo limita la rigidez torsional, y los elementos que giran los limita la fatiga.
Dimensionar todo contra el límite elástico produce un coche pesado que además
falla por donde no se miraba.

\subsection{Tensiones básicas}

\begin{align}
  \sigma_{\text{axial}} &= \frac{N}{A} &
  \sigma_{\text{flexión}} &= \frac{M\,y}{I} &
  \tau_{\text{torsión}}  &= \frac{T\,r}{J}
  \label{eq:tensiones_basicas}
\end{align}

Cuando actúan varias a la vez, se combinan con el criterio de von Mises, que es
el que se compara con el límite elástico:

\begin{equation}
  \sigma_{VM} = \sqrt{\sigma_x^2 - \sigma_x\sigma_y + \sigma_y^2 + 3\tau_{xy}^2}
  \label{eq:von_mises}
\end{equation}

\subsection{Pandeo: lo que gobierna los trapecios}

Una barra esbelta a compresión no falla por aplastamiento: colapsa lateralmente
mucho antes. La carga crítica de Euler es

\begin{equation}
  P_{cr} = \frac{\pi^2 E I}{(K L)^2}
  \label{eq:euler}
\end{equation}

donde $L$ es la longitud entre apoyos y $K$ depende de cómo estén sujetos los
extremos:

\begin{center}
\begin{tabular}{@{}lcl@{}}
\toprule
\textbf{Condición de extremos} & $K$ & \textbf{Caso típico} \\
\midrule
Articulado--articulado & 1{,}0 & \textbf{Barra de suspensión con rótulas} \\
Empotrado--articulado  & 0{,}7 & Extremo soldado, otro rotulado \\
Empotrado--empotrado   & 0{,}5 & Barra soldada en ambos extremos \\
Empotrado--libre       & 2{,}0 & Voladizo \\
\bottomrule
\end{tabular}
\end{center}

Para el caso normal en suspensión --- rótula en los dos extremos --- $K = 1$, que
es el peor caso razonable. Conviene usarlo aunque haya algo de empotramiento: la
diferencia entre $K=1$ y $K=0{,}7$ duplica la carga crítica, y esa no es una
apuesta que merezca la pena hacer.

\begin{clave}
Para un tubo de pared delgada de radio medio $r_m$ y espesor $t$:
$A \approx 2\pi r_m t$ e $I \approx \pi r_m^3 t$. Sustituyendo en
\eqref{eq:euler} con la \emph{masa constante} (es decir, $r_m t$ fijo), resulta
$P_{cr} \propto r_m^2$.

Traducido: \textbf{duplicar el diámetro de un tubo, adelgazando la pared para
mantener el peso, multiplica por cuatro la carga de pandeo.} Por eso los
trapecios son tubos de diámetro grande y pared fina, y por eso el límite práctico
no lo pone el cálculo sino la pared mínima que se puede soldar, el pandeo local
de la pared y lo que exija el reglamento.
\end{clave}

La fórmula de Euler solo vale para barras esbeltas. Para esbelteces bajas
($\lambda = KL/r_g$ pequeño, con $r_g=\sqrt{I/A}$ el radio de giro) el fallo pasa
a ser plastificación y hay que usar una curva de transición --- por ejemplo la
parábola de Johnson. En la práctica: si la carga de Euler sale por encima de la
carga de plastificación $A\,S_y$, el pandeo no es el modo crítico y hay que
comprobar resistencia.

% ---------------------------------------------------------------------
\section{Fatiga}
\label{sec:mecanica_aplicada:03}
\index{fatiga}

Una pieza sometida a carga cíclica puede romper con tensiones muy por debajo del
límite elástico. La grieta nace en la superficie, casi siempre en un cambio de
sección, y crece hasta que la sección restante no aguanta.

\subsection{¿Importa la fatiga en un coche de Formula Student?}

Depende de la pieza, y merece pensarlo antes de gastar tiempo:

\begin{center}
\begin{tabularx}{\textwidth}{@{}lXl@{}}
\toprule
\textbf{Tipo de pieza} & \textbf{Por qué} & \textbf{¿Fatiga?} \\
\midrule
Palieres, bujes, tornillos de rueda & Giran: miles de ciclos por kilómetro & \textbf{Sí, siempre} \\
Anclajes del chasis, trapecios & Carga variable, muchos km de test & Sí \\
Uniones soldadas & El cordón concentra tensión & Sí \\
Piezas que se reutilizan otra temporada & Acumulan historia desconocida & Sí \\
Estructura cargada casi estáticamente & Pocos ciclos de amplitud alta & Normalmente no \\
\bottomrule
\end{tabularx}
\end{center}

La resistencia de \SI{22}{\kilo\meter} induce a pensar que la fatiga no importa.
Es engañoso por tres motivos: la temporada de ensayos son cientos de kilómetros,
las piezas que giran ven un ciclo completamente invertido por vuelta de rueda, y
las piezas heredadas de temporadas anteriores llegan con un daño acumulado que
nadie ha contabilizado.

\subsection{Límite de fatiga}

Para aceros, el límite de fatiga en probeta se estima como
$S'_e \approx 0{,}5\,S_{ut}$ (con un techo en torno a \SI{700}{\mega\pascal}). Ese
valor de laboratorio se corrige con los factores de Marin para llegar al de la
pieza real:

\begin{equation}
  S_e = k_a\,k_b\,k_c\,k_d\,k_e\; S'_e
  \label{eq:marin}
\end{equation}

donde $k_a$ recoge el acabado superficial, $k_b$ el tamaño, $k_c$ el tipo de
carga, $k_d$ la temperatura y $k_e$ la fiabilidad exigida. El más castigador
suele ser $k_a$: una superficie bruta de laminación puede reducir el límite de
fatiga a menos de la mitad respecto de una pulida.

\begin{ojo}
El aluminio \textbf{no tiene límite de fatiga}: su curva S--N sigue bajando
indefinidamente. Una pieza de aluminio no se puede diseñar «para vida infinita»;
hay que fijar un número de ciclos objetivo y una vida en kilómetros, y retirarla
o inspeccionarla al alcanzarla. Esto afecta directamente a manguetas, balancines
y anclajes mecanizados.
\end{ojo}

\subsection{Concentración de tensiones y tensión media}

El efecto de una entalla se recoge con el factor de concentración $K_t$
(geométrico) corregido por la sensibilidad a la entalla $q$ del material:

\begin{equation}
  K_f = 1 + q\,(K_t - 1)
  \label{eq:kf}
\end{equation}

Y cuando la carga cíclica se superpone a una carga media --- lo habitual --- se
comprueba con un criterio como el de Goodman:

\begin{equation}
  \frac{\sigma_a}{S_e} + \frac{\sigma_m}{S_{ut}} = \frac{1}{n}
  \label{eq:goodman}
\end{equation}

\begin{clave}
En fatiga, la geometría manda sobre el material. Un radio de acuerdo generoso, un
buen acabado superficial y evitar cambios bruscos de sección hacen más por la
vida de una pieza que subir de grado de material --- que además, al ser más
resistente, suele ser \emph{más} sensible a la entalla.
\end{clave}

% ---------------------------------------------------------------------
\section{Uniones atornilladas y pares de apriete}
\label{sec:mecanica_aplicada:04}
\index{uniones atornilladas}\index{par de apriete}

Buena parte de los fallos de un coche de competición aparecen en las uniones
antes que en las piezas, y muchos de ellos se explican por una precarga
incorrecta.

\subsection{La precarga lo es todo}

Un tornillo bien apretado no es un tornillo que sujeta: es un tornillo que
\emph{comprime} las piezas entre sí con una fuerza mucho mayor que la carga
externa. Mientras esa compresión no desaparezca, la unión no se mueve, no vibra y
--- esto es lo importante --- el tornillo apenas ve la carga variable.

La precarga recomendada es una fracción alta de la carga de prueba del tornillo:

\begin{equation}
  F_i = 0{,}75\, A_t\, S_p \quad\text{(uniones desmontables)}, \qquad
  F_i = 0{,}90\, A_t\, S_p \quad\text{(uniones permanentes)}
  \label{eq:precarga}
\end{equation}

con $A_t$ el área resistente a tracción y $S_p$ la carga de prueba del grado:

\begin{center}
\begin{tabular}{@{}lSS@{}}
\toprule
\textbf{Grado} & {$S_p$ (\si{\mega\pascal})} & {$S_{ut}$ (\si{\mega\pascal})} \\
\midrule
8.8  & 600 & 800 \\
10.9 & 830 & 1040 \\
12.9 & 970 & 1220 \\
\bottomrule
\end{tabular}
\end{center}

\subsection{Del par de apriete a la precarga}

La relación práctica es

\begin{equation}
  T = K\, d\, F_i
  \label{eq:par_apriete}
\end{equation}

con $d$ el diámetro nominal y $K$ un coeficiente que vale aproximadamente
\num{0.20} en seco y \num{0.15} lubricado.

\begin{ojo}
El coeficiente $K$ tiene una dispersión enorme --- fácilmente $\pm\SI{25}{\percent}$
--- porque depende del rozamiento en la rosca y bajo la cabeza, que a su vez
depende del acabado, la lubricación y si el tornillo es nuevo. Es decir:
\textbf{una llave dinamométrica controla el par, no la precarga}. Es lo mejor que
tenemos en el taller, pero hay que saber que el número real de precarga tiene esa
incertidumbre encima, y no diseñar uniones que dependan de acertarla con precisión.
\end{ojo}

\subsection{Reparto de la carga externa}

La rigidez relativa del tornillo $k_b$ y de las piezas $k_m$ determina qué parte
de una carga externa $P$ ve el tornillo:

\begin{equation}
  C = \frac{k_b}{k_b + k_m}, \qquad
  F_{\text{tornillo}} = F_i + C\,P, \qquad
  F_{\text{unión}} = F_i - (1-C)\,P
  \label{eq:reparto_carga}
\end{equation}

Como las piezas suelen ser mucho más rígidas que el tornillo, $C$ es pequeño
(típicamente \numrange{0.1}{0.3}) y el tornillo apenas nota la carga variable. De
ahí una regla de diseño poco intuitiva: \textbf{un tornillo largo (menos rígido)
aguanta mejor a fatiga que uno corto}, porque reduce $C$.

La unión se separa cuando $F_i - (1-C)P = 0$. Diseñar con margen frente a ese
punto es más importante que comprobar la tensión del tornillo.

\subsection{Reglas prácticas}

\begin{itemize}
  \item \textbf{Cortante: por aplastamiento, no por rozamiento.} Salvo que la
        precarga esté controlada de verdad, se dimensiona suponiendo que el
        tornillo trabaja a cortante contra el agujero. Y la rosca nunca debe
        quedar en el plano de cortadura: se usa tornillo de vástago liso.
  \item \textbf{Longitud de rosca engranada}: al menos $1\,d$ en acero y
        $2\,d$ en aluminio. En aluminio, mejor inserto roscado.
  \item \textbf{Arandelas} siempre bajo la cabeza que gira, y bajo cualquier
        apoyo en aluminio.
  \item \textbf{No reutilizar} tornillos apretados cerca del límite elástico, ni
        tuercas autoblocantes de anillo plástico que ya hayan trabajado.
  \item \textbf{Marcar el apriete}: una raya de pintura entre tornillo y pieza
        permite ver de un vistazo si algo se ha aflojado, tanto en el taller como
        en la revisión técnica.
\end{itemize}

\begin{regla}[frenado de tornillería]
El reglamento exige \emph{frenado positivo} en los tornillos críticos: elementos
que impidan mecánicamente el aflojamiento, como tuercas autoblocantes, tuerca
almenada con pasador o alambre de frenado. Los adhesivos de bloqueo de roscas
habitualmente \textbf{no} se aceptan por sí solos como frenado positivo.
Compruébalo en el artículo correspondiente del reglamento vigente antes de
diseñar la unión: es un motivo clásico de rechazo en la revisión técnica.
\end{regla}

% ---------------------------------------------------------------------
\section{Selección de materiales}
\label{sec:mecanica_aplicada:05}
\index{selección de materiales}

\subsection{El dato que casi nadie espera}

La rigidez específica --- módulo elástico dividido por densidad --- de casi todos
los metales estructurales es prácticamente la misma:

\begin{center}
\begin{tabular}{@{}lSSS@{}}
\toprule
\textbf{Material} & {$E$ (\si{\giga\pascal})} & {$\rho$ (\si{\gram\per\cubic\centi\meter})} & {$E/\rho$} \\
\midrule
Acero (cualquier aleación) & 210 & 7.85 & 26.8 \\
Aluminio                   & 70  & 2.70 & 25.9 \\
Titanio                    & 110 & 4.51 & 24.4 \\
Magnesio                   & 45  & 1.74 & 25.9 \\
\bottomrule
\end{tabular}
\end{center}

\begin{clave}
Dos consecuencias que ahorran muchas discusiones:

\begin{enumerate}
  \item \textbf{Un acero aleado no es más rígido que un acero normal.} Todos los
        aceros tienen $E \approx \SI{210}{\giga\pascal}$. Se elige un acero
        aleado por resistencia, no por rigidez.
  \item \textbf{A igualdad de masa, y si la pieza trabaja a tracción pura, da
        igual el metal.} La ventaja del aluminio aparece cuando la pieza trabaja a
        \emph{flexión o pandeo}: como es menos denso, con la misma masa se hace
        más voluminosa, y el momento de inercia crece mucho más deprisa que la
        pérdida de módulo. Es un efecto de \emph{forma}, no de material.
\end{enumerate}
\end{clave}

Donde sí hay diferencias enormes es en la resistencia específica, y ahí es donde
la elección de aleación y tratamiento importa de verdad.

\subsection{Paleta habitual en Formula Student}

\begin{center}
\begin{tabularx}{\textwidth}{@{}llX@{}}
\toprule
\textbf{Familia} & \textbf{Designación} & \textbf{Cuándo y avisos} \\
\midrule
Acero al carbono & S235, S355 & Utillaje, soportes poco cargados. Barato y disponible \\
Acero aleado & 25CrMo4 & Tubos de chasis y piezas soldadas muy cargadas. Soldable \\
Acero aleado & 42CrMo4 & Ejes y palieres. Requiere tratamiento térmico \\
Aluminio & 6082-T6 & Uso general mecanizado. \textbf{Soldable}, buena corrosión \\
Aluminio & 7075-T6 & Manguetas, balancines. Alta resistencia. \textbf{No soldable}, mala corrosión \\
Titanio & Ti-6Al-4V & Solo si está justificado con números. Caro y difícil de mecanizar \\
\bottomrule
\end{tabularx}
\end{center}

\pendiente{Ajustar esta tabla a lo que el equipo puede comprar y mecanizar de
verdad, con proveedores y plazos. Un material que no se puede conseguir en dos
semanas en el formato necesario no es una opción de diseño.}

\subsection{Método de selección}

\begin{enumerate}
  \item \textbf{Identifica el modo limitante} (resistencia, rigidez, pandeo,
        fatiga) y la geometría disponible. Este paso decide casi todo.
  \item \textbf{Escribe el índice de material} que corresponde a ese modo
        (\textcite{ashby_materials_2016} los tabula todos) y ordena candidatos.
  \item \textbf{Filtra por fabricabilidad}: ¿se puede soldar, mecanizar, tratar
        con lo que tenemos o con lo que nos pueden hacer?
  \item \textbf{Filtra por disponibilidad y coste}: formatos comerciales, plazo de
        entrega, y lo que costará en el informe del capítulo
        \ref{ch:evento_costes}.
  \item \textbf{Comprueba compatibilidad}: contacto galvánico entre metales
        distintos, temperatura de servicio, y compatibilidad con adhesivos o
        fluidos (capítulo \ref{ch:fabricacion_ii}).
\end{enumerate}

\begin{ojo}
\begin{itemize}
  \item Elegir el material antes que la forma. El material es la última decisión,
        no la primera.
  \item Usar 7075 y descubrir en el taller que hay que soldarlo.
  \item Especificar un tratamiento térmico y no comprobar quién lo va a hacer,
        cuánto tarda y cuánto deforma la pieza.
  \item Poner aluminio y acero en contacto directo en una zona que se moja.
  \item Copiar el material de la pieza equivalente de otro equipo sin saber qué
        modo de fallo estaban combatiendo.
\end{itemize}
\end{ojo}

% ---------------------------------------------------------------------
%  Cierre del capítulo
% ---------------------------------------------------------------------

\begin{caso}
\redactar{Rellenar con nuestros casos reales: qué se ha roto en el coche y por
qué. Para cada fallo, indicar el modo (estático, pandeo, fatiga, unión floja), si
se había calculado y con qué caso de carga. Es el recuadro más valioso del
capítulo y solo se puede escribir mirando piezas rotas.}
\end{caso}

\begin{ojo}
\begin{itemize}
  \item Dimensionar todo contra el límite elástico sin preguntarse qué limita
        realmente la pieza.
  \item Aplicar un coeficiente de seguridad genérico a un caso de carga que nadie
        ha definido: el coeficiente no arregla que la carga esté mal.
  \item Apretar «a sensación». La diferencia entre un tornillo M8 bien apretado y
        uno apretado con fuerza es de más del doble de precarga.
  \item No documentar los pares de apriete: si no están en una tabla, en el
        montaje cada uno usa el suyo.
\end{itemize}
\end{ojo}

\begin{entregable}
\begin{enumerate}
  \item Una hoja de cálculo de \textbf{pandeo de tubos}: entradas diámetro,
        espesor, longitud y condiciones de extremo; salida carga crítica y
        coeficiente de seguridad.
  \item Una hoja de cálculo de \textbf{uniones atornilladas}: entradas grado,
        métrica y tipo de unión; salidas precarga, par de apriete y carga de
        separación.
  \item La \textbf{tabla de pares de apriete del coche}, por tipo de unión,
        impresa y colgada en el taller.
\end{enumerate}
\end{entregable}

\begin{juez}
\begin{enumerate}
  \item ¿Qué limita el dimensionado de este trapecio: resistencia, rigidez o pandeo?
  \item ¿Por qué este tubo tiene este diámetro y este espesor y no otros?
  \item ¿Con qué par se aprieta esto y de dónde sale ese número?
  \item ¿Has comprobado fatiga en alguna pieza? ¿En cuáles y por qué en esas?
  \item ¿Por qué esta pieza es de 7075 y no de 6082?
\end{enumerate}
\end{juez}

% !TeX root = ../../main.tex
% =====================================================================
%  Análisis por elementos finitos con criterio
%  Responsable:
%  Última revisión:
% =====================================================================
\chapter{Análisis por elementos finitos con criterio}
\label{ch:elementos_finitos}

\begin{objetivos}
  \item Decidir si un problema merece un FEA o se resuelve más rápido y mejor a
        mano.
  \item Reconocer una singularidad numérica y no informar nunca de la tensión que
        aparece en ella.
  \item Justificar las condiciones de contorno de un análisis, que es de donde
        salen la mayoría de los errores.
  \item Hacer una comprobación manual que confirme o desmienta el resultado.
  \item Redactar un informe de análisis que otra persona pueda revisar.
\end{objetivos}

\fichasesion{90 minutos}{Capítulo \ref{ch:mecanica_aplicada}}%
{Ordenador con el software de simulación y una pieza real del coche}

Un programa de elementos finitos siempre da un resultado, y el resultado siempre
tiene el mismo aspecto convincente tanto si el modelo está bien planteado como si
no. La habilidad que hay que adquirir, por tanto, no es hacer correr el análisis,
sino \emph{saber cuándo el resultado no vale}. \textcite{cook_fea_2001} es la
referencia de fondo si se quiere entender lo que hay debajo.

% ---------------------------------------------------------------------
\section{Qué puede y qué no puede decirte un FEA}
\label{sec:elementos_finitos:01}

\begin{clave}
Un FEA no analiza tu pieza: analiza \textbf{el modelo que has construido de tu
pieza}. Si la geometría está simplificada, si las cargas son inventadas o si los
apoyos no se parecen a la realidad, el resultado es exacto \emph{para un problema
que no es el tuyo}.
\end{clave}

\begin{center}
\begin{tabularx}{\textwidth}{@{}XX@{}}
\toprule
\textbf{Bien, y de forma fiable} & \textbf{Mal, o solo con mucho cuidado} \\
\midrule
Rigidez y desplazamientos & Tensión máxima absoluta en una entalla \\
Comparar dos diseños entre sí & Vida a fatiga \\
Ver por dónde van las cargas & Uniones atornilladas y contactos \\
Decidir dónde quitar material & Cordones de soldadura \\
Modos y frecuencias propias & Pandeo, si solo se hace estático lineal \\
Reacciones en los apoyos & Materiales compuestos \\
\bottomrule
\end{tabularx}
\end{center}

Dos consecuencias prácticas. La primera: para \textbf{comparar} dos diseños, el
FEA es excelente aunque el modelo tenga defectos, siempre que los defectos sean
los mismos en ambos casos. La segunda: para dar un \textbf{número absoluto} que
va a justificar un coeficiente de seguridad, el listón es mucho más alto y hay
que verificar y validar.

Y una regla previa a todo lo demás: \textbf{si no sabes aproximadamente qué va a
salir, no lances el análisis}. Estima antes el orden de magnitud a mano; esa
estimación es lo que después te permitirá detectar que el resultado es absurdo.

% ---------------------------------------------------------------------
\section{Mallado y convergencia}
\label{sec:elementos_finitos:02}
\index{mallado}\index{convergencia}

\subsection{Elegir el tipo de elemento}

\begin{itemize}
  \item \textbf{Nunca uses tetraedros de primer orden para calcular tensiones.}
        Son excesivamente rígidos y dan resultados sistemáticamente bajos. Si el
        programa ofrece «borrador» o «rápido», normalmente es esto.
  \item \textbf{Tetraedros de segundo orden} son la opción por defecto razonable
        para piezas macizas.
  \item \textbf{Elementos lámina} para todo lo que sea de pared delgada: chapas,
        laminados, tubos si interesa el detalle local.
  \item \textbf{Elementos viga} para estructuras tubulares. Un análisis de rigidez
        torsional de un chasis tubular (capítulo \ref{ch:rigidez_torsional}) se
        hace con vigas: es más rápido, más limpio de interpretar y más fácil de
        comprobar a mano que un modelo sólido con cientos de miles de elementos.
\end{itemize}

\subsection{El estudio de convergencia}

Un resultado sin estudio de convergencia no es un resultado. El procedimiento es
mecánico:

\begin{enumerate}
  \item Elige la magnitud que te interesa: un desplazamiento, una tensión en un
        punto concreto, una reacción.
  \item Resuelve con tres o cuatro tamaños de malla decrecientes.
  \item Representa esa magnitud frente al número de grados de libertad.
  \item Si la curva se aplana, has convergido: quédate con la malla a partir de
        la cual el cambio es menor que un umbral que tú fijas (por ejemplo
        \SI{2}{\percent}).
  \item \textbf{Si la curva no se aplana y sigue creciendo, no es un problema de
        malla: es una singularidad.}
\end{enumerate}

\subsection{Singularidades}
\index{singularidad}

En una arista viva entrante, en una carga puntual o en un apoyo puntual, la
tensión teórica es infinita. Al refinar la malla, la tensión que informa el
programa crece sin parar y nunca converge. No es un fallo del programa: la
solución exacta del \emph{modelo} es infinita ahí.

\begin{ojo}
Informar de la tensión máxima que aparece en una arista viva es un error que un
juez detecta nada más ver la imagen. La respuesta correcta a «¿cuál es la tensión
máxima?» empieza por «esa zona es una singularidad; el valor que doy es el de
aquí, a esta distancia, y esta es la razón».
\end{ojo}

Qué hacer cuando aparece una:

\begin{itemize}
  \item \textbf{Modelar el radio de acuerdo real.} Casi siempre la pieza no tiene
        una arista viva: la tiene el modelo, porque alguien la simplificó.
  \item \textbf{Repartir la carga o el apoyo} en una superficie en lugar de en un
        punto o una arista.
  \item \textbf{Evaluar la tensión a una distancia} de la singularidad, indicando
        cuál, si el detalle local no es lo que se está estudiando.
  \item \textbf{Cambiar de magnitud}: si lo que importa es la rigidez, usa el
        desplazamiento, que converge mucho antes y no tiene este problema.
\end{itemize}

% ---------------------------------------------------------------------
\section{Condiciones de contorno}
\label{sec:elementos_finitos:03}
\index{condiciones de contorno}

Aquí es donde se pierden los análisis. La geometría suele estar bien y el
material también; lo que casi nunca se parece a la realidad es cómo se ha sujetado
y cargado la pieza.

\subsection{Sobrerrestricción}

Empotrar una cara entera porque «ahí va atornillada» añade una rigidez que no
existe y genera tensiones enormes justo en el borde del empotramiento, que además
son falsas. La pieza real está sujeta por unos tornillos que permiten cierto giro
y cierto deslizamiento.

\begin{center}
\begin{tabularx}{\textwidth}{@{}XX@{}}
\toprule
\textbf{En lugar de} & \textbf{Haz} \\
\midrule
Empotrar la cara de apoyo completa & Restringir solo lo necesario para evitar
movimiento de sólido rígido \\
Empotrar el agujero de un tornillo & Un punto remoto acoplado al agujero, o una
carga tipo apoyo distribuida \\
Fijar un nodo para «quitar el sólido rígido» & Alivio inercial, o tres
restricciones mínimas bien colocadas \\
Contacto pegado en toda la unión & Modelar el contacto real, o justificar por qué
pegado es conservador \\
\bottomrule
\end{tabularx}
\end{center}

\subsection{Simetría}

Solo se puede usar simetría cuando \textbf{son simétricas las tres cosas}:
geometría, cargas y apoyos. Un caso de carga en curva no es simétrico. Un caso de
frenada en recta sí puede serlo. Aprovechar la simetría cuando no la hay elimina
del modelo justo la parte de la carga que descompensa la pieza.

\subsection{Las cargas vienen de algún sitio}

Las fuerzas de un análisis no se eligen: salen de los casos de carga definidos en
el capítulo \ref{ch:casos_carga}, que a su vez salen de la dinámica del vehículo.
Un análisis con cargas «razonables» inventadas por quien lo hace no demuestra
nada, y en el evento de diseño la primera pregunta después de ver la imagen
siempre es de dónde sale ese número.

% ---------------------------------------------------------------------
\section{Verificación con cálculo manual}
\label{sec:elementos_finitos:04}
\index{verificación}

Antes de creerte un resultado, tres comprobaciones que cuestan poco:

\begin{enumerate}
  \item \textbf{Equilibrio.} La suma de las reacciones tiene que dar la suma de
        las cargas aplicadas, en las tres direcciones. Si no cuadra, hay una carga
        o un apoyo donde no debía.
  \item \textbf{Orden de magnitud a mano.} Simplifica la pieza a una viga, un
        tubo o una placa y aplica las fórmulas del capítulo
        \ref{ch:mecanica_aplicada}. Si el FEA y la mano difieren en un factor 2,
        puede ser la geometría; si difieren en un factor 10, hay un error.
  \item \textbf{Un caso conocido con la misma configuración.} Resuelve una viga
        en voladizo con tu mismo tipo de elemento, tu mismo mallado y tu mismo
        tipo de apoyo, y compárala con la solución analítica. Si eso no sale, lo
        demás tampoco.
\end{enumerate}

\begin{clave}
Conviene distinguir dos cosas que se confunden constantemente:
\textbf{verificación} es comprobar que estás resolviendo bien las ecuaciones
(malla, convergencia, equilibrio); \textbf{validación} es comprobar que las
ecuaciones que resuelves describen la realidad, y eso solo se hace midiendo la
pieza real. El banco de rigidez torsional del capítulo
\ref{ch:rigidez_torsional} es validación; un estudio de convergencia es
verificación. Hacen falta las dos.
\end{clave}

% ---------------------------------------------------------------------
\section{Cómo se documenta un análisis}
\label{sec:elementos_finitos:05}

Un análisis que no está documentado no se puede revisar, no se puede defender
ante un juez y no sirve el año que viene. El formato es una página, y siempre la
misma:

\begin{center}
\begin{tabularx}{\textwidth}{@{}lX@{}}
\toprule
\textbf{Apartado} & \textbf{Qué contiene} \\
\midrule
Objetivo & Qué decisión de diseño depende de este análisis \\
Geometría & Modelo usado y \emph{qué se ha simplificado y por qué} \\
Material & Designación y de dónde salen las propiedades \\
Caso de carga & Valores y referencia al documento del que salen \\
Contorno & Apoyos y cargas, con su justificación física \\
Malla & Tipo de elemento, tamaño y \textbf{estudio de convergencia} \\
Resultados & Magnitudes de interés, con leyenda y unidades \\
Comprobación & El cálculo manual y su comparación \\
Conclusión & Coeficiente de seguridad y \textbf{qué se cambia en el diseño} \\
\bottomrule
\end{tabularx}
\end{center}

Sobre las imágenes: llevan leyenda con escala y unidades, la escala no puede estar
autoajustada al valor de una singularidad, y si el desplazamiento está amplificado
hay que decir por cuánto. Sin esos tres datos, la imagen no se puede interpretar.

% ---------------------------------------------------------------------
%  Cierre del capítulo
% ---------------------------------------------------------------------

\begin{caso}[con qué calculamos]
El equipo está aprendiendo \textbf{Ansys}, y hasta ahora los análisis se han
hecho con los módulos de simulación integrados en \textbf{Fusion 360} y
\textbf{SolidWorks}. Conviene saber qué implica cada opción:

\begin{itemize}
  \item Los módulos integrados en el CAD son cómodos porque la geometría ya está
        ahí, pero dan poco control sobre el tipo de elemento, el mallado y los
        apoyos, y aplican valores por defecto que no siempre se ven. Sirven para
        comparar alternativas; cuesta más justificar con ellos un número absoluto.
  \item Ansys da ese control, y es la razón por la que merece la pena el esfuerzo
        de aprenderlo: permite elegir el orden del elemento y hacer de verdad el
        estudio de convergencia del apartado
        \ref{sec:elementos_finitos:02}.
\end{itemize}

Dos reglas mientras conviven las tres herramientas:

\begin{enumerate}
  \item \textbf{Cada persona resuelve el caso de referencia en la herramienta que
        vaya a usar} antes de analizar nada del coche. Es la única forma de saber
        qué hacen los valores por defecto de tu programa.
  \item \textbf{No se comparan resultados obtenidos con programas distintos} para
        decidir entre dos diseños: estarías comparando las herramientas, no los
        diseños. Una comparación se hace siempre con el mismo montaje.
\end{enumerate}

\redactar{Añadir aquí, cuando exista, un análisis real completo siguiendo la
plantilla del apartado \ref{sec:elementos_finitos:05}. Lo ideal es uno en el que
el resultado obligara a cambiar el diseño, y mejor todavía si se pudo contrastar
con una pieza rota o con una medida.}
\end{caso}

\begin{ojo}
\begin{itemize}
  \item Informar de la tensión máxima en una arista viva.
  \item Decir «ha convergido» refiriéndose a que el solver terminó, no a que la
        malla esté convergida. Son cosas distintas.
  \item Comparar dos diseños con mallas diferentes: la diferencia que ves puede
        ser solo la malla.
  \item Dar un coeficiente de seguridad contra un caso de carga que nadie ha
        definido por escrito.
  \item Comprobar el límite elástico en una pieza cuyo problema real es la
        rigidez o el pandeo.
  \item Mallar más fino para «afinar el resultado» sin comprobar si lo que hay
        debajo es una singularidad.
  \item Enseñar la imagen deformada con factor \num{200} sin decirlo.
\end{itemize}
\end{ojo}

\begin{entregable}
\begin{enumerate}
  \item La \textbf{plantilla de informe de análisis} del equipo, con los nueve
        apartados de la tabla anterior.
  \item Un \textbf{caso de referencia resuelto} (viga en voladizo) con la
        configuración habitual del equipo, contrastado con la solución analítica,
        guardado para que cualquiera pueda comprobar su montaje.
  \item El \textbf{checklist de revisión de un FEA}, que usa el revisor antes de
        dar por bueno un análisis ajeno.
\end{enumerate}
\end{entregable}

\begin{juez}
\begin{enumerate}
  \item ¿De dónde salen las cargas de este análisis?
  \item ¿Cómo has sujetado la pieza en el modelo y por qué se parece eso a la
        realidad?
  \item ¿Has hecho estudio de convergencia? Enséñamelo.
  \item Esa tensión máxima está en una arista viva. ¿Qué significa ese número?
  \item ¿Con qué has comprobado que este resultado es correcto?
  \item ¿Qué has cambiado en el diseño después de este análisis?
\end{enumerate}
\end{juez}

% !TeX root = ../../main.tex
% =====================================================================
%  Ensayo y adquisición de datos
%  Responsable: 
%  Última revisión: 
% =====================================================================
\chapter{Ensayo y adquisición de datos}
\label{ch:ensayo_adquisicion_datos}

\begin{clave}
\textbf{Capítulo pendiente, y a propósito.} El equipo todavía no tiene campaña de
ensayos ni sistema de adquisición de datos, y este libro no describe métodos que
no practicamos: sería justo el tipo de contenido copiado que hace inútil un
manual interno.

Se redactará cuando haya datos propios que contar. Mientras tanto, los capítulos
que dependen de él ---\ref{ch:piloto_ensayos_pista} (ensayos en pista),
\ref{ch:cfd_validacion} (validación de CFD) y \ref{ch:validacion_termica}
(validación térmica)--- quedan igualmente incompletos, y conviene tenerlo
presente: hoy el equipo diseña con modelos que nadie ha contrastado con la
realidad. Cerrar este capítulo es el paso que convierte el resto del libro en
ingeniería verificada.
\end{clave}

\pendiente{Redactar cuando exista la primera campaña de ensayos. El esqueleto de
apartados de abajo se conserva como guion de partida.}

\section{Qué medir y por qué}
\label{sec:ensayo_adquisicion_datos:01}
\pendiente{Pendiente de redactar.}

\section{Sensores habituales en el coche}
\label{sec:ensayo_adquisicion_datos:02}
\pendiente{Pendiente de redactar.}

\section{Muestreo, filtrado y ruido}
\label{sec:ensayo_adquisicion_datos:03}
\pendiente{Pendiente de redactar.}

\section{Diseño de un ensayo}
\label{sec:ensayo_adquisicion_datos:04}
\pendiente{Pendiente de redactar.}

\section{Del dato a la decisión}
\label{sec:ensayo_adquisicion_datos:05}
\pendiente{Pendiente de redactar.}

% ---------------------------------------------------------------------
%  Cierre del capítulo: los cuatro recuadros son obligatorios.
% ---------------------------------------------------------------------

\begin{caso}
\redactar{Qué hicimos nosotros: decisión tomada, datos de partida, resultado en pista y qué haríamos distinto.}
\end{caso}

\begin{ojo}
\begin{itemize}
  \item \redactar{Error que repetimos cada temporada en este tema.}
\end{itemize}
\end{ojo}

\begin{entregable}
\redactar{Qué produce el lector al terminar: plantilla, hoja de cálculo, checklist o apartado del informe.}
\end{entregable}

\begin{juez}
\begin{enumerate}
  \item \redactar{Pregunta tipo del design event sobre este capítulo.}
\end{enumerate}
\end{juez}



% =====================================================================
\backmatter
% =====================================================================
\printbibliography[heading=bibintoc,title={Bibliografía}]

\cleardoublepage
\addcontentsline{toc}{chapter}{Lo que ya sabemos que falta}
\listoftodos[Lo que ya sabemos que falta]

\end{document}
